\chapter{线性系统综合：状态反馈、观测器与分离原理}
\label{ch:linear-synthesis}

% ════════════════════════════════════════════════════════════════
%  P4 控制部「C0 现代控制理论基础」子部 · 第 C0-4 章
%  主源：刘豹《现代控制理论》(第3版) 第5章（状态反馈/输出反馈基本结构·极点配置
%        定理与证明·系统镇定·解耦[前馈/状态反馈/能解耦判据/标准形]·状态观测器
%        [定义/存在性/全维设计/微分器放噪动机链]·降维观测器·带观测器状态反馈+分离原理）
%  辅源/第二信源：张嗣瀛/高立群《现代控制理论》(第2版) 第6章 6.1–6.5（Ackermann/Bass-Gura
%        算法·反馈解耦 E 阵非奇异充要条件·降维观测器）+ ★6.6 鲁棒入门一节
%        （Schur 补/LMI/H∞ 状态反馈 ARE/回旋倒立摆与两足机器人实例）——02 独有人审入口。
%  与导论 ch:control 的密度梯度去重：导论 sec:ctrl-pole-placement 已"给结论会用"
%        （极点配置定理、Luenberger eq:ctrl-luenberger、误差动态 eq:ctrl-obs-error、
%        分离原理 prop:ctrl-separation），本章"讲透+推一遍+建直觉"，严格化交 C/E/F 章。
%  OCR 勘误已逐式核对：刘豹 (A^T+c^T G^T)^T=A+Gc（去导数点，M2-b L2322）；解耦 d_i 范围
%        n-1（M2-b L2659）；例 b=[1;0]（M2-c L3747）；降维观测器去多余负号（M2-c L3332）。
%        现控02：里卡蒂(6.6.6) PA+A^T P-PBB^T P+Q_0=0（BB^T 非 B^T B，L3506）；增益 K=-B^T P
%        （去衍生 x，L3509）；γ^2（L3538）；H∞ ARE 补闭括号后 P（L4029）。
% ════════════════════════════════════════════════════════════════

\begin{practice}[前置自测]
开始之前，先试答下列问题。若已能流畅作答，可只速读本章的框图与速查卡；若卡住，正是本章要为你解决的。
\begin{enumerate}
  \item 给一个可控的线性系统，你想让它的闭环极点恰好落在指定位置（比如临界阻尼、某个带宽）。状态反馈 $\mathbf u=-\mathbf K\mathbf x$ 凭什么能做到？“可控”在这里扮演什么角色？
  \item 状态反馈要用到\emph{全部}状态 $\mathbf x$，可机器人上往往只有编码器测得位置、测不到所有速度与内部量。一个朴素的想法是“把位置微分一下得速度”——为什么工程上不能这么干？
  \item 既然测不到 $\mathbf x$，就造一个“假的 $\mathbf x$”——用系统模型在线跑一个副本 $\hat{\mathbf x}$。但模型副本会漂移，怎么让它\emph{自动纠偏}、稳稳跟上真实状态？
  \item 把估计的 $\hat{\mathbf x}$ 喂给状态反馈 $\mathbf u=-\mathbf K\hat{\mathbf x}$，控制器的 $\mathbf K$ 与观测器的 $\mathbf L$ 会不会“互相打架”？能不能各设各的、拼起来还稳？
  \item 多输入多输出系统里，动一个输入会同时晃动好几个输出（耦合）。能否用反馈把它“拆”成一组互不干扰的单回路？什么样的系统拆得开、什么样的拆不开？
  \item 真实模型总有误差（$\mathbf A,\mathbf B$ 不准）。精确极点配置依赖精确模型——当模型本身带一团不确定性 $\Delta\mathbf A,\Delta\mathbf B$ 时，还能保证镇定吗？
\end{enumerate}
\end{practice}

\section*{本章目标}
学完本章，你应当能够：
\begin{enumerate}
  \item \textbf{说清}状态反馈、输出反馈、从输出到 $\dot{\mathbf x}$ 反馈三种基本反馈结构，并论证状态反馈为何比输出反馈“自由度更大、性能更优”；
  \item \textbf{独立证明}极点配置定理（状态反馈可任意配置闭环极点的充要条件是系统完全可控），并掌握两条等价的设计路线——经可控标准型的系数比对法与 Ackermann / Bass--Gura 公式；
  \item \textbf{辨析}镇定与极点配置的关系，给出“状态反馈可镇定 $\iff$ 不可控子系统渐近稳定”这一充要条件及其结构分解证明；
  \item \textbf{设计}多变量系统的解耦控制——理解前馈补偿器解耦与状态反馈解耦的区别，掌握特征量 $d_i,\mathbf E_i$ 与“可解耦 $\iff$ $\mathbf E$ 非奇异”这一判据，并知道积分型解耦系统为何还需追加极点配置；
  \item \textbf{建立}状态观测器的动机链——从“状态不可测”出发，理解“微分器重构状态”为何因放大噪声而不可行，进而引出“模型副本 + 输出残差纠正”的渐近观测器；
  \item \textbf{推导}全维（Luenberger）观测器的误差动态 $\dot{\tilde{\mathbf x}}=(\mathbf A-\mathbf L\mathbf C)\tilde{\mathbf x}$，用对偶原理把观测器增益 $\mathbf L$ 的设计化为一个极点配置问题，并掌握观测器存在的充要条件；
  \item \textbf{构造}降维观测器：利用 $\mathbf y$ 直接提供 $m$ 个状态分量、只对其余 $n-m$ 个未测分量造 $(n-m)$ 维观测器；
  \item \textbf{运用}分离原理：证明带观测器状态反馈的闭环极点恰为控制器极点与观测器极点之并，从而 $\mathbf K$ 与 $\mathbf L$ 可独立设计，理解它是“先估计、再控制”工程范式的理论许可证；
  \item \textbf{入门}线性不确定系统的鲁棒镇定——用 Schur 补与 Lyapunov 方法把镇定条件写成线性矩阵不等式（LMI），理解 $H_\infty$ 状态反馈的代数 Riccati 方程，并能在回旋倒立摆 / 两足机器人模型上看到它的落地。
\end{enumerate}

\subsection*{与导论的分工}
本章是 P4 控制部「现代控制理论基础」子部（C0）的第四章，紧接 \cref{ch:lyapunov-basics}（稳定性与李雅普诺夫方法）之后。\textbf{它与导论 \cref{ch:control} 的关系是“同一主题、不同密度”}：导论 \cref{sec:ctrl-pole-placement} 已把极点配置定理、Ackermann 公式、Luenberger 观测器（\cref{eq:ctrl-luenberger}）、观测器误差动态（\cref{eq:ctrl-obs-error}）、分离原理（\cref{prop:ctrl-separation}）以\emph{导论密度}陈述清楚——“给结论、会用”。本章不重复这些编号，而是\textbf{把它们讲透：先建物理直觉与动机链（状态不可测$\to$微分器放噪$\to$龙伯格$\to$分离原理），再把极点配置定理推一遍、把观测器从存在性到全维 / 降维设计走通，最后铺到解耦与鲁棒入门}。一句话：导论让你\emph{会用}状态反馈与观测器，本章让你\emph{会推、会设计、且理解为什么}。凡导论已立的 def / thm，本章开篇半句交叉引用接住即下沉展开（如“展开 \cref{prop:ctrl-separation}”），决不另起炉灶；而严格证明与前沿（分离原理的确定性等价证明、$H_\infty$ 博弈型 Riccati、CBF 安全对偶）则交由 A--Q 高级章。

\subsection*{本章知识导航}
本章是一条“\emph{从知道全部状态的理想反馈，一步步走到只用可测输出、且对模型误差鲁棒的可实现反馈}”的主线。结构如下：

\begin{center}
\small
\begin{tikzpicture}[
  node distance=4mm and 7mm, >=Stealth,
  blk/.style={draw, rounded corners, align=center, font=\small, inner sep=3pt, fill=main!6, draw=main!60!black},
  grp/.style={draw, rounded corners, align=center, font=\small, inner sep=3pt, fill=second!6, draw=second!60!black},
  red/.style={draw, rounded corners, align=center, font=\small, inner sep=3pt, fill=third!8, draw=third!60!black},
  ln/.style={->, thick, gray!60!black}]
\node[blk] (fb) {反馈结构\\状态/输出/到$\dot{\mathbf x}$};
\node[blk, right=of fb] (pole) {极点配置\\可控$\Rightarrow$任意配};
\node[blk, right=of pole] (stab) {镇定\\不可控子系统稳};
\node[grp, below=8mm of fb] (decoup) {解耦\\$\mathbf E$ 非奇异};
\node[grp, right=of decoup] (obs) {全维观测器\\模型副本+纠错};
\node[grp, right=of obs] (red) {降维观测器\\$\mathbf y$ 直供 $m$ 维};
\node[red, below=8mm of decoup] (sep) {分离原理\\$\mathbf K,\mathbf L$ 独立};
\node[red, right=of sep] (lqg) {LQG / 估计-控制对偶};
\node[red, right=of lqg] (robust) {鲁棒入门\\Schur/LMI/$H_\infty$};
\draw[ln] (fb)--(pole); \draw[ln] (pole)--(stab);
\draw[ln] (stab.south) -- ++(0,-4mm) -| (decoup.north);
\draw[ln] (decoup)--(obs); \draw[ln] (obs)--(red);
\draw[ln] (red.south) -- ++(0,-4mm) -| (sep.north);
\draw[ln] (sep)--(lqg); \draw[ln] (lqg)--(robust);
\end{tikzpicture}
\end{center}

\noindent\textbf{推荐路径}：\cref{sec:syn-feedback,sec:syn-pole} 是主干（反馈结构 + 极点配置，建立“可控 $\Rightarrow$ 可任意配极点”这一核心定理）；\cref{sec:syn-stabilize} 的镇定是极点配置的弱化版（只要求稳定、不要求精确配置），面向“原系统不稳、先救活”的场景；\cref{sec:syn-decoupling}（解耦）面向多变量系统，可按需选读；\cref{sec:syn-observer,sec:syn-reduced} 是另一条主干（观测器：从动机链到全维 / 降维设计），与极点配置经\emph{对偶}精确镜像；\cref{sec:syn-separation}（分离原理）把控制器与观测器焊接，是全章的会合点，也是“先估计、再控制”范式的理论核心；\cref{sec:syn-robust}（鲁棒入门）是面向模型不确定性的一节人审入口，用 LMI / $H_\infty$ 语言给出在倒立摆 / 两足机器人上的落地，深推交 \cref{ch:robust-hinf}。

\subsection*{前置知识桥接}
本章坐落在 C0 子部前三章的地基之上：状态空间模型 $\dot{\mathbf x}=\mathbf A\mathbf x+\mathbf B\mathbf u,\ \mathbf y=\mathbf C\mathbf x$ 与可控 / 可观标准型（\cref{ch:statespace-basics,ch:controllability-obsv}），可控性是极点配置的前提、可观性是观测器存在的前提（\cref{ch:controllability-obsv}），Lyapunov 第二法与李氏方程 $\mathbf A^\top\mathbf P+\mathbf P\mathbf A=-\mathbf Q$（镇定与鲁棒镇定的工具，\cref{ch:lyapunov-basics}）。在数学侧，本章用到 P1 的半正定矩阵与 LMI / 凸优化（\cref{ch:nlopt}，鲁棒一节把镇定写成 LMI），以及 P0 的机械系统动力学记号 $\mathbf M(\mathbf q)\ddot{\mathbf q}+\mathbf C\dot{\mathbf q}+\mathbf g=\boldsymbol\tau$（\cref{ch:rigid_body}，倒立摆实例的母体）。横向上，观测器与 \cref{ch:eskf} 的卡尔曼滤波“结构同形、增益来路不同”，本章在分离原理处与之显式互引。

\begin{table}[H]\centering\small
\caption{本章阅读方式建议}
\begin{tabular}{lll}
\toprule
阅读方式 & 时间 & 适合谁 \\
\midrule
精读（含全部推导与例题） & 4--6 小时 & 首次系统学习、要做反馈 / 观测器设计 \\
速读（跳过冗长矩阵计算） & 1.5 小时 & 学过线性系统、想对齐本书记号与设计流程 \\
速查（只看小结的符号表 + 设计流程速查表） & 15 分钟 & 设计控制器 / 观测器时回来查步骤 \\
\bottomrule
\end{tabular}
\end{table}

\begin{note}
\textbf{本章记号与约定.}\;
状态 $\mathbf x\in\mathbb R^n$、输入 $\mathbf u\in\mathbb R^r$、输出 $\mathbf y\in\mathbb R^m$ 一律粗体（两本教材源用非粗体 / $\boldsymbol x$，融入时一律粗体化）。系统三元组记 $\Sigma_0=(\mathbf A,\mathbf B,\mathbf C)$；单输入 / 单输出时记 $(\mathbf A,\mathbf b,\mathbf c)$。
\textbf{状态反馈增益}统一记 $\mathbf K$（闭环阵 $\mathbf A-\mathbf B\mathbf K$，反馈律 $\mathbf u=\mathbf v-\mathbf K\mathbf x$，$\mathbf v$ 为参考输入；教材源有时记 $\mathbf u=\mathbf v+\mathbf K\mathbf x$ 而吸收负号，本书统一取\emph{负反馈} $-\mathbf K$、与导论 \cref{sec:ctrl-pole-placement} 和 LQR \cref{sec:ctrl-lqr} 一致）。
\textbf{观测器增益}统一记 $\mathbf L$（误差阵 $\mathbf A-\mathbf L\mathbf C$，与导论 \cref{eq:ctrl-luenberger} 一致）；两本教材源观测器增益记 $\mathbf G$ 或 $\mathbf E$，融入时一律改 $\mathbf L$。
解耦的输入变换阵记 $\mathbf F$（$\mathbf u=\mathbf F\mathbf v-\mathbf K\mathbf x$）。
李氏 / Riccati 矩阵记 $\mathbf P\succ\mathbf 0$，李氏方程统一 $\mathbf A^\top\mathbf P+\mathbf P\mathbf A=-\mathbf Q$（与 \cref{eq:lyap-eq} 同形）。
\textbf{$\mathbf P/\mathbf Q/\mathbf K/\mathbf L$ 双关纪律}：本章 $\mathbf K$ 是\emph{控制}增益、$\mathbf L$ 是\emph{观测器}增益；勿与 \cref{ch:eskf} 估计侧的 Kalman 增益（那里记 $\mathbf K_{\rm Kalman}$ 或 $\mathbf L$）混淆——两者经“估计--控制对偶”相关，但语义不同。
\textbf{“可控 / 可观”与“能控 / 能观”在本书同义}（教材多用“能控 / 能观”，本书正文用“可控 / 可观”对齐导论）。
\end{note}

% ════════════════════════════════════════════════════════════════
\section{反馈控制的基本结构}
\label{sec:syn-feedback}

\paragraph{动机：综合就是“给系统装一个控制器”。}
前三章是\emph{分析}——给定系统，判别它可控、可观、稳定与否。本章转入\emph{综合（synthesis）}：给定一个受控对象 $\Sigma_0=(\mathbf A,\mathbf B,\mathbf C)$ 和一组性能指标（稳定、指定极点、解耦、抗扰），设计一个反馈控制器，使闭环系统满足这些指标\cite{liubao2006modern}。现代控制理论的综合与经典控制理论一样，仍是“受控对象 + 反馈控制器”的闭环结构；区别在于：经典理论习惯用\emph{输出反馈}，现代理论更多用\emph{状态反馈}——因为状态承载了系统的全部内部信息，反馈它能提供更丰富的自由度。

\paragraph{三种基本反馈。}
设受控对象
\begin{equation}
  \dot{\mathbf x}=\mathbf A\mathbf x+\mathbf B\mathbf u,\qquad \mathbf y=\mathbf C\mathbf x,
  \label{eq:syn-plant}
\end{equation}
按反馈信号取自何处，有三种基本形式\cite{liubao2006modern}：

\begin{enumerate}
  \item \textbf{状态反馈}：把每个状态分量乘以增益反馈到输入端，$\mathbf u=\mathbf v-\mathbf K\mathbf x$（$\mathbf K$ 为 $r\times n$ 增益阵，$\mathbf v$ 为参考输入）。代入 \cref{eq:syn-plant}：
  \begin{equation}
    \dot{\mathbf x}=(\mathbf A-\mathbf B\mathbf K)\mathbf x+\mathbf B\mathbf v,\qquad \mathbf y=\mathbf C\mathbf x.
    \label{eq:syn-state-fb}
  \end{equation}
  闭环系统记 $\Sigma_K=((\mathbf A-\mathbf B\mathbf K),\,\mathbf B,\,\mathbf C)$\footnote{源（刘豹 §5.1.1）此处闭环简记的括号被 OCR 弄丢一层（M2-b L1770），本书据上下文补全为 $((\cdot),\mathbf B,\mathbf C)$。},维数不变，但 $\mathbf K$ 的选择可\emph{自由改变闭环特征值}。
  \item \textbf{输出反馈}：把输出乘以增益反馈，$\mathbf u=\mathbf v-\mathbf H\mathbf y=\mathbf v-\mathbf H\mathbf C\mathbf x$，闭环 $\dot{\mathbf x}=(\mathbf A-\mathbf B\mathbf H\mathbf C)\mathbf x+\mathbf B\mathbf v$。
  \item \textbf{从输出到 $\dot{\mathbf x}$ 反馈}：把输出反馈到状态导数处，$\dot{\mathbf x}=(\mathbf A-\mathbf G\mathbf C)\mathbf x+\mathbf B\mathbf u$——这种形式正是\emph{状态观测器}的数学骨架（\cref{sec:syn-observer}），故单独列出。
\end{enumerate}

\begin{figure}[H]
\centering
\begin{tikzpicture}[>=Stealth, node distance=8mm, font=\small,
  box/.style={draw, minimum height=7mm, minimum width=14mm, fill=main!6},
  sum/.style={draw, circle, inner sep=1pt, minimum size=4mm}]
  \node (v) at (0,0) {$\mathbf v$};
  \node[sum, right=8mm of v] (s1) {$+$};
  \node[box, right=10mm of s1] (plant) {$\dot{\mathbf x}=\mathbf A\mathbf x+\mathbf B\mathbf u$};
  \node[right=14mm of plant] (y) {$\mathbf y=\mathbf C\mathbf x$};
  \node[box, below=9mm of plant] (K) {$\mathbf K$};
  \draw[->] (v)--(s1);
  \draw[->] (s1)--(plant) node[midway,above]{$\mathbf u$};
  \draw[->] (plant)--(y);
  \coordinate (tap) at ($(plant.east)+(7mm,0)$);
  \draw[->] (tap) |- (K) node[pos=0.25,right]{$\mathbf x$};
  \draw[->] (K) -| (s1) node[pos=0.85,left]{$-$};
\end{tikzpicture}
\caption{多输入多输出系统的状态反馈结构：状态 $\mathbf x$ 经增益阵 $\mathbf K$ 负反馈到输入端，与参考 $\mathbf v$ 相加构成控制律 $\mathbf u=\mathbf v-\mathbf K\mathbf x$。（仿刘豹 §5.1 图 5.1 重绘）}
\label{fig:syn-state-fb}
\end{figure}

\begin{insight}
\textbf{为什么状态反馈优于输出反馈？}\;
比较两式：输出反馈中的 $\mathbf H\mathbf C$ 恰相当于状态反馈中的 $\mathbf K$。但 $\mathbf y\in\mathbb R^m$ 而 $\mathbf x\in\mathbb R^n$，通常 $m<n$，故 $\mathbf H$（$r\times m$）可调元素远少于 $\mathbf K$（$r\times n$）——输出反馈只相当于一种\emph{部分}状态反馈，自由度受限。\textbf{边界}：仅当 $\mathbf C=\mathbf I$（输出即全状态）时 $\mathbf H\mathbf C=\mathbf K$，二者等同。代价是状态反馈需要\emph{全部}状态可得——这正是后面观测器要解决的难题。输出反馈的突出优点则是\emph{技术实现方便}（输出本就可测，无需重构）。
\end{insight}

\paragraph{反馈对可控 / 可观性的影响。}
一个关键事实：\textbf{状态反馈不改变可控性，但可能破坏可观性}\cite{liubao2006modern,zhang2017modern}。可控性不变是因为 $\Sigma_0$ 与 $\Sigma_K$ 的可控性矩阵满足
\begin{equation}
  [\,\mathbf B\ \ (\mathbf A-\mathbf B\mathbf K)\mathbf B\ \cdots\ ]=[\,\mathbf B\ \ \mathbf A\mathbf B\ \cdots\ ]\begin{bmatrix}\mathbf I&*&\cdots\\&\mathbf I&\\&&\ddots\end{bmatrix},
  \label{eq:syn-ctrb-invariant}
\end{equation}
右乘的是非奇异（块上三角、对角为 $\mathbf I$）阵，秩不变。可观性可能被破坏，是因为状态反馈\emph{改变极点但不改变零点}：闭环传递函数分子（零点）不变、分母（极点）随 $\mathbf K$ 改变，于是有可能制造出\emph{零极点对消}，对消掉的那个模态便不再可观（\cref{sec:co-cancel} 的零极点对消正是此现象的源头）。对偶地，输出反馈不改变可观性、但可能破坏可控性。

% ════════════════════════════════════════════════════════════════
\section{极点配置}
\label{sec:syn-pole}

\paragraph{动机：把性能指标翻译成极点位置。}
线性系统的动态品质（超调、调节时间、阻尼）主要由极点在复平面的分布决定。综合的一种最直接形式便是：给定一组\emph{期望极点} $\{\lambda_i^*\}$（由时域指标换算而来），通过反馈把闭环极点恰好配到这些位置。这就是\textbf{极点配置（pole placement）}问题。它把“我想要怎样的响应”这一含糊愿望，化为“把极点放到 $\{\lambda_i^*\}$”这一精确的代数任务。

\begin{theorem}[极点配置定理]
\label{thm:syn-pole-placement}
对单输入系统 $\Sigma_0=(\mathbf A,\mathbf b,\mathbf c)$，采用状态反馈 $\mathbf u=\mathbf v-\mathbf K\mathbf x$ 可将闭环 $(\mathbf A-\mathbf b\mathbf K)$ 的极点配置到复平面上\emph{任意指定}的（关于实轴对称的）$n$ 个位置，其\textbf{充要条件}是 $\Sigma_0$ 完全可控\cite{liubao2006modern,wonham1967pole}。多输入系统结论相同（充要条件仍是完全可控），但 $\mathbf K$ 的解不唯一、设计更复杂。
\end{theorem}

\noindent{}（这是导论 \cref{sec:ctrl-pole-placement} 所陈述结论的展开版；导论给“可控 $\Rightarrow$ 任意配”，本章补\emph{充要性的构造性证明}与\emph{设计算法}。）

\begin{proof}[充分性证明（单输入，经可控标准型）]
设 $\Sigma_0$ 完全可控。证明分四步\cite{liubao2006modern}。

\emph{第一步：化可控标准 I 型。}\;可控保证存在非奇异变换 $\mathbf x=\mathbf T_c\bar{\mathbf x}$，把 $\Sigma_0$ 化为可控标准 I 型（友矩阵形，\cref{sec:co-canonical}）：
\begin{equation}
  \bar{\mathbf A}=\mathbf T_c^{-1}\mathbf A\mathbf T_c=\begin{bmatrix}0&1&\cdots&0\\\vdots&&\ddots&\vdots\\0&0&\cdots&1\\-a_0&-a_1&\cdots&-a_{n-1}\end{bmatrix},\quad \bar{\mathbf b}=\mathbf T_c^{-1}\mathbf b=\begin{bmatrix}0\\\vdots\\0\\1\end{bmatrix},
  \label{eq:syn-ccf}
\end{equation}
其末行 $\{-a_i\}$ 恰是开环特征多项式 $\det(s\mathbf I-\mathbf A)=s^n+a_{n-1}s^{n-1}+\cdots+a_1 s+a_0$ 的系数。

\emph{第二步：在标准型坐标下加反馈。}\;取 $\bar{\mathbf K}=[\bar k_0\ \bar k_1\ \cdots\ \bar k_{n-1}]$，则闭环阵 $\bar{\mathbf A}-\bar{\mathbf b}\bar{\mathbf K}$ 只有\emph{末行}被改动：
\begin{equation}
  \bar{\mathbf A}-\bar{\mathbf b}\bar{\mathbf K}=\begin{bmatrix}0&1&\cdots&0\\\vdots&&\ddots&\vdots\\0&0&\cdots&1\\-(a_0+\bar k_0)&-(a_1+\bar k_1)&\cdots&-(a_{n-1}+\bar k_{n-1})\end{bmatrix},
  \label{eq:syn-ccf-closed}
\end{equation}
仍是友矩阵形，故闭环特征多项式可\emph{直接读出}：
\begin{equation}
  f(\lambda)=\lambda^n+(a_{n-1}+\bar k_{n-1})\lambda^{n-1}+\cdots+(a_1+\bar k_1)\lambda+(a_0+\bar k_0).
  \label{eq:syn-closed-char}
\end{equation}

\emph{第三步：系数比对。}\;令 $f(\lambda)$ 等于期望多项式 $f^*(\lambda)=\prod_{i=1}^n(\lambda-\lambda_i^*)=\lambda^n+a^*_{n-1}\lambda^{n-1}+\cdots+a^*_0$，逐项比对同次幂系数：
\begin{equation}
  \boxed{\bar k_i=a^*_i-a_i,\qquad i=0,1,\dots,n-1.}
  \label{eq:syn-kbar}
\end{equation}
友矩阵的末行可\emph{独立、无约束}地配到任意值，故 $f(\lambda)=f^*(\lambda)$ 总能实现——这就是“可控 $\Rightarrow$ 可任意配”的代数本质。

\emph{第四步：变回原坐标。}\;由 $\mathbf u=\mathbf v-\bar{\mathbf K}\bar{\mathbf x}=\mathbf v-\bar{\mathbf K}\mathbf T_c^{-1}\mathbf x$，得原坐标增益
\begin{equation}
  \mathbf K=\bar{\mathbf K}\,\mathbf T_c^{-1}.
  \label{eq:syn-K-orig}
\end{equation}
充分性证毕。
\end{proof}

\begin{insight}
\textbf{必要性的直觉（为何不可控就配不动）。}\;若 $\Sigma_0$ 不可控，由 Kalman 结构分解（\cref{sec:co-decomp}）存在坐标使 $\mathbf A=\big[\begin{smallmatrix}\mathbf A_c&\mathbf A_{12}\\\mathbf 0&\mathbf A_{\bar c}\end{smallmatrix}\big],\ \mathbf b=\big[\begin{smallmatrix}\mathbf b_c\\\mathbf 0\end{smallmatrix}\big]$。任何状态反馈 $\mathbf K=[\mathbf K_c\ \mathbf K_{\bar c}]$ 作用后，闭环阵的\emph{右下块仍是 $\mathbf A_{\bar c}$}（反馈进不到不可控子系统）——$\mathbf A_{\bar c}$ 的特征值\textbf{雷打不动}，配不动。可控部分的极点可任意配，不可控部分的极点完全锁死。这就是必要性：要\emph{任意}配全部 $n$ 个极点，必须没有锁死的部分，即完全可控。
\end{insight}

\begin{note}
\textbf{多输入系统：充要条件相同，但 $\mathbf K$ 不唯一.}\;
对多输入系统（$\mathbf B$ 为 $n\times r,\ r>1$），极点配置的充要条件仍是完全可控，但设计更复杂、且 $\mathbf K$（$r\times n$）的解\emph{不唯一}\cite{liubao2006modern,zhang2017modern}。原因是：$n$ 个极点约束远少于 $\mathbf K$ 的 $r\times n$ 个自由度，故配同一组极点有无穷多个 $\mathbf K$。这种\emph{冗余自由度}是好事——可用来\emph{兼顾}极点配置以外的指标（如同时解耦、最小化增益范数、配置特征矢量方向）。但化多输入系统为标准型颇繁，工程上常转而用数值算法（如 Tits--Yang 的鲁棒极点配置，把多余自由度用于最大化对扰动的鲁棒性），或干脆改用 LQR——后者用代价权重 $\mathbf Q,\mathbf R$ \emph{隐式}选定一组“最优”极点，避开了“极点放哪、多余自由度怎么用”的人工抉择（\cref{ch:optimal-control-basics}）。
\end{note}

\begin{algo}[多输入极点配置：龙伯格标准形法（张嗣瀛\cite{zhang2017modern} §6.2.4）]
\label{algo:syn-multiinput}
设 $(\mathbf A,\mathbf B)$ 完全可控、$\operatorname{rank}\mathbf B=r$，求 $\mathbf K$（$r\times n$）使 $\operatorname{eig}(\mathbf A-\mathbf B\mathbf K)=\{\lambda_i^*\}$。给出单输入“化标准型$\to$比系数”三步的多输入推广。
\begin{enumerate}
  \item \textbf{定可控性指数 $\mu_i$.}\;在可控性矩阵 $\mathcal C=[\,\mathbf b_1,\dots,\mathbf b_r,\ \mathbf A\mathbf b_1,\dots,\mathbf A\mathbf b_r,\ \dots,\ \mathbf A^{n-1}\mathbf b_1,\dots,\mathbf A^{n-1}\mathbf b_r\,]$ 中\emph{从左到右}挑出 $n$ 个线性无关列，按
  \[
    [\,\mathbf b_1,\mathbf A\mathbf b_1,\dots,\mathbf A^{\mu_1-1}\mathbf b_1,\ \mathbf b_2,\dots,\mathbf A^{\mu_2-1}\mathbf b_2,\ \dots,\ \mathbf b_r,\dots,\mathbf A^{\mu_r-1}\mathbf b_r\,]=:\mathbf Q^{-1}
  \]
  重排，由此读出\emph{可控性指数} $\mu_1,\dots,\mu_r$（$\sum_{i=1}^r\mu_i=n$）。
  \item \textbf{构造变换阵 $\mathbf P$.}\;取 $\mathbf Q$ 各\emph{块末行} $\mathbf e_{1\mu_1}^\top,\dots,\mathbf e_{r\mu_r}^\top$，叠成
  \[
    \mathbf P=\big[\,\mathbf e_{1\mu_1}^\top;\ \mathbf e_{1\mu_1}^\top\mathbf A;\ \dots;\ \mathbf e_{1\mu_1}^\top\mathbf A^{\mu_1-1};\ \mathbf e_{2\mu_2}^\top;\ \dots;\ \mathbf e_{r\mu_r}^\top\mathbf A^{\mu_r-1}\,\big],
  \]
  经 $\tilde{\mathbf x}=\mathbf P\mathbf x$ 化为\textbf{龙伯格标准形} $(\tilde{\mathbf A},\tilde{\mathbf B})$：$\tilde{\mathbf A}$ 分块，对角块 $\tilde{\mathbf A}_{ii}$ 是 $\mu_i\times\mu_i$ 友矩阵，$\tilde{\mathbf B}$ 仅在各块末行有单位元（友矩阵的“可独立改写的末行”由此从单输入的一行变成 $r$ 行）。
  \item \textbf{分块配置.}\;在标准形坐标取 $\tilde{\mathbf K}$（$r\times n$）使 $\tilde{\mathbf A}-\tilde{\mathbf B}\tilde{\mathbf K}$ 成\emph{单一友矩阵}（整体特征多项式 $=\prod_i(\lambda-\lambda_i^*)$），或成\emph{块对角友矩阵} $\operatorname{diag}(\tilde{\mathbf A}_{11},\dots,\tilde{\mathbf A}_{rr})$（各块独立分担一组期望极点，$\prod_i$ 各块特征多项式 $=\prod_i(\lambda-\lambda_i^*)$），由系数比对定 $\tilde{\mathbf K}$。
  \item \textbf{变回原坐标.}\;$\mathbf K=\tilde{\mathbf K}\mathbf P$。
\end{enumerate}
\end{algo}

\begin{note}
\textbf{多输入“$\mathbf K$ 不唯一”的算法层面具体化.}\;\cref{algo:syn-multiinput} 第 3 步的“单一友矩阵”与“块对角”两种选法，对\emph{同一组}期望极点给出\emph{不同}的 $\mathbf K$——这正是 \cref{thm:syn-pole-placement} 后注“多输入 $\mathbf K$ 不唯一”的算法体现。如张嗣瀛\cite{zhang2017modern} 的一个五阶双输入例（期望极点 $-1,-2\pm\mathrm j,-1\pm2\mathrm j$），两种选法分别得 $\mathbf K=\big[\begin{smallmatrix}2&0&0&0&1\\25&55&48&23&5\end{smallmatrix}\big]$ 与 $\mathbf K=\big[\begin{smallmatrix}7&9&5&1&1\\0&0&0&1&0\end{smallmatrix}\big]$，二者都使 $(\mathbf A-\mathbf B\mathbf K)$ 取得同一组特征值。这一冗余自由度，正是上文“可挪作鲁棒化 / 解耦 / 最小增益”之用的来源；工程上更常直接调数值例程（\texttt{place} / 鲁棒极点配置）或用 LQR 隐式选极点，规避手工化龙伯格标准形之繁。
\end{note}

\begin{example}[状态反馈配置极点]
\label{ex:syn-pole}
设系统传递函数 $W(s)=\dfrac{10}{s(s+1)(s+2)}$，求状态反馈使闭环极点为 $-2,\,-1\pm j$\cite{liubao2006modern}。

\textbf{解.}\;传递函数无零极点对消 $\Rightarrow$ 系统可控可观。直接写可控标准 I 型实现：$a_0=0,a_1=2,a_2=3$。期望多项式 $f^*(\lambda)=(\lambda+2)(\lambda+1+j)(\lambda+1-j)=\lambda^3+4\lambda^2+6\lambda+4$，故 $a^*_0=4,a^*_1=6,a^*_2=4$。由 \cref{eq:syn-kbar}：$\bar{\mathbf K}=[\,a^*_0-a_0,\ a^*_1-a_1,\ a^*_2-a_2\,]=[\,4,\ 4,\ 1\,]$（标准型坐标）。再经 \cref{eq:syn-K-orig} 变回原坐标即得 $\mathbf K$。\emph{注意}：若一开始就用可控标准型，可免去坐标变换、直接由系数算 $\bar{\mathbf K}$；但标准型的状态分量往往\emph{难以直接检测}，工程实现反需借助观测器（\cref{sec:syn-observer}）或按物理可测变量选坐标。
\end{example}

\subsection{Ackermann 公式与 Bass--Gura 算法}
\label{sec:syn-ackermann}

经可控标准型的“化型$\to$比系数$\to$变回”三步，可压缩成一个\emph{闭式}。导论 \cref{sec:ctrl-pole-placement} 已给 \textbf{Ackermann 公式}\cite{ackermann1972der}：单输入可控系统的反馈增益为
\begin{equation}
  \mathbf K=[\,0\ \cdots\ 0\ \ 1\,]\,\mathcal C^{-1}\,f^*(\mathbf A),
  \label{eq:syn-ackermann}
\end{equation}
其中 $\mathcal C=[\,\mathbf b\ \ \mathbf A\mathbf b\ \cdots\ \mathbf A^{n-1}\mathbf b\,]$ 是可控性矩阵（\cref{eq:ctrl-cmat}），$f^*(\mathbf A)=\mathbf A^n+a^*_{n-1}\mathbf A^{n-1}+\cdots+a^*_0\mathbf I$ 是把期望多项式代入矩阵 $\mathbf A$。可控性保证 $\mathcal C$ 可逆。\cref{eq:syn-ackermann} 的妙处在于：无需显式化标准型，一步算出增益。

\begin{insight}
\textbf{Ackermann 公式为何长这样？}\;关键是 Cayley--Hamilton 定理（\cref{ch:controllability-obsv}）：闭环阵满足\emph{它自己}的特征方程，即 $f^*(\mathbf A-\mathbf b\mathbf K)=\mathbf 0$。把 $f^*(\mathbf A-\mathbf b\mathbf K)$ 按 $\mathbf b\mathbf K$ 展开并与 $f^*(\mathbf A)$ 相减，所有含 $\mathbf K$ 的项都左带一个 $\mathbf b$ 因子、且按 $\mathbf b,\mathbf A\mathbf b,\dots,\mathbf A^{n-1}\mathbf b$ 排列——于是 $f^*(\mathbf A)=\mathcal C\cdot(\text{含}\mathbf K\text{的列})$。左乘 $\mathcal C^{-1}$ 解出该列，其\emph{末行}恰为所求增益（前 $n-1$ 行是中间量），故用行选择子 $[0\cdots0\ 1]$ 取出，即 \cref{eq:syn-ackermann}。一句话：Ackermann 把“化标准型 + 比系数”三步，用 Cayley--Hamilton 一次性\emph{代数地}完成。
\end{insight}

现控理论02 给出一条等价的\textbf{Bass--Gura 型算法}\cite{zhang2017modern}，对手算更友好（设 $\det(s\mathbf I-\mathbf A)=s^n+a_1 s^{n-1}+\cdots+a_n$，期望 $\alpha^*(s)=s^n+a^*_1 s^{n-1}+\cdots+a^*_n$）：

\begin{algo}[Bass--Gura 单输入极点配置]
\label{algo:syn-bassgura}
\begin{enumerate}
  \item 求开环特征多项式系数 $a_1,\dots,a_n$ 与期望多项式系数 $a^*_1,\dots,a^*_n$；
  \item 令 $\tilde{\mathbf K}=[\,a^*_n-a_n\ \ a^*_{n-1}-a_{n-1}\ \cdots\ a^*_1-a_1\,]$；
  \item 构造 $\mathbf Q=[\,\mathbf b\ \ \mathbf A\mathbf b\ \cdots\ \mathbf A^{n-1}\mathbf b\,]\cdot\mathbf W$，其中 $\mathbf W$ 是由 $\{a_i\}$ 构成的上三角 Toeplitz 阵（首副对角为 $1$）；
  \item 令 $\mathbf P=\mathbf Q^{-1}$，则 $\mathbf K=\tilde{\mathbf K}\mathbf P$。
\end{enumerate}
\end{algo}

\begin{example}[Bass--Gura 算例]
\label{ex:syn-bassgura}
$\dot{\mathbf x}=\big[\begin{smallmatrix}0&0&0\\1&-1&0\\0&1&-1\end{smallmatrix}\big]\mathbf x+\big[\begin{smallmatrix}1\\0\\0\end{smallmatrix}\big]\mathbf u$，配 $\lambda^*=-2,\,-1\pm j\sqrt3$\cite{zhang2017modern}。可控（$\operatorname{rank}\mathcal C=3$）。开环 $\det(s\mathbf I-\mathbf A)=s^3+2s^2+s$，得 $a_1=2,a_2=1,a_3=0$；期望 $(s+2)(s+1+j\sqrt3)(s+1-j\sqrt3)=s^3+4s^2+8s+8$，得 $a^*_1=4,a^*_2=8,a^*_3=8$。于是 $\tilde{\mathbf K}=[\,a^*_3-a_3,\ a^*_2-a_2,\ a^*_1-a_1\,]=[\,8,\ 7,\ 2\,]$，再乘 $\mathbf P=\mathbf Q^{-1}$ 即得 $\mathbf K$。
\end{example}

\begin{pitfall}[把“可控”当成“能任意快地驱动”]
极点配置理论上能把极点配到任意远（任意快），但\textbf{可控性只保证“能配”，不保证“代价有界”}。配得越快（极点越往左），所需控制量越大、对模型误差越敏感、越易触发执行器饱和。正确理解：可控性是\emph{定性}的“能否”，速度 / 代价是\emph{定量}问题——由 LQR 的权重 $\mathbf Q,\mathbf R$ 折中（\cref{ch:optimal-control-basics}），而非一味追求最左极点。这与\emph{观测器}极点配得太快会放大噪声（\cref{sec:syn-observer} 的动机链）是同一枚硬币的两面。
\end{pitfall}

\subsection{输出反馈与从输出到 \texorpdfstring{$\dot{\mathbf x}$}{x-dot} 反馈配置}
\label{sec:syn-output-pole}

\paragraph{输出反馈：自由度受限的极点配置。}
若坚持只用\emph{输出反馈} $\mathbf u=\mathbf v-\mathbf H\mathbf y$（不重构状态），闭环阵为 $\mathbf A-\mathbf B\mathbf H\mathbf C$。由于 $\mathbf H$ 的可调元素只有 $r\times m$ 个（远少于状态反馈的 $r\times n$），\emph{一般无法任意配置全部 $n$ 个极点}——可达的闭环极点只是参数空间里一张低维曲面。这是 \cref{sec:syn-feedback} 那条“输出反馈 = 部分状态反馈”论断在极点配置层面的具体后果。要突破这一限制而又只用可测的 $\mathbf y$，唯一出路便是\emph{先重构状态}（观测器）再做状态反馈——分离原理（\cref{thm:syn-separation}）保证这条路覆盖全空间。

\paragraph{从输出到 $\dot{\mathbf x}$ 反馈：观测器极点配置的对偶骨架。}
第三种反馈 $\dot{\mathbf x}=(\mathbf A-\mathbf G\mathbf C)\mathbf x+\mathbf B\mathbf u$（输出反馈到状态导数）虽不是物理控制律，却是\emph{观测器设计}的代数骨架。其极点配置问题“选 $\mathbf G$ 配 $(\mathbf A-\mathbf G\mathbf C)$ 的极点”，经转置 $(\mathbf A-\mathbf G\mathbf C)^\top=\mathbf A^\top-\mathbf C^\top\mathbf G^\top$ 与对偶原理（\cref{thm:co-duality}），\textbf{完全等价于对对偶系统 $(\mathbf A^\top,\mathbf C^\top)$ 设计状态反馈 $\mathbf G^\top$}\cite{liubao2006modern}。于是“配 $\mathbf A-\mathbf G\mathbf C$ 的极点”可任意实现 $\iff$ $(\mathbf A^\top,\mathbf C^\top)$ 完全可控 $\iff$ $(\mathbf A,\mathbf C)$ 完全可观。这正是 \cref{sec:syn-observer} 观测器增益设计的预演——把它单列于此，是为了让“状态反馈配极点”与“观测器配误差极点”这对孪生问题在同一处对照。

\begin{insight}
\textbf{一张表看懂三种反馈的孪生结构。}\;状态反馈配 $(\mathbf A-\mathbf B\mathbf K)$、要 $(\mathbf A,\mathbf B)$ 可控；从输出到 $\dot{\mathbf x}$ 反馈配 $(\mathbf A-\mathbf G\mathbf C)$、要 $(\mathbf A,\mathbf C)$ 可观——二者经 $\mathbf A\leftrightarrow\mathbf A^\top,\ \mathbf B\leftrightarrow\mathbf C^\top,\ \mathbf K\leftrightarrow\mathbf G^\top$ 精确镜像。\textbf{记住这组对偶替换，本章后半（观测器、降维、分离）的每个结论都可由前半（极点配置、镇定）“免费”翻译得到}——这是把综合理论从“两套公式”压成“一套 + 对偶”的关键认知。
\end{insight}

% ════════════════════════════════════════════════════════════════
\section{系统镇定}
\label{sec:syn-stabilize}

\paragraph{动机：有时只要“活下来”，不必精确配极点。}
保证稳定是控制系统正常工作的前提。\textbf{镇定（stabilization）}问题问的是：对受控对象 $\Sigma_0$，能否通过反馈使其闭环极点全部具负实部（渐近稳定）？这是极点配置的\emph{弱化}版——不要求把极点精确放到指定位置，只要求统统挪到左半平面。于是，即便系统不可控（无法任意配极点），仍可能可镇定：只需把那些\emph{不稳定}的模态搬进左半平面即可。

\begin{theorem}[状态反馈可镇定的充要条件]
\label{thm:syn-stabilize}
对系统 $\Sigma_0=(\mathbf A,\mathbf B,\mathbf C)$，采用状态反馈可镇定的充要条件是其\textbf{不可控子系统渐近稳定}\cite{liubao2006modern,zhang2017modern}。
\end{theorem}

\begin{proof}[证明（经结构分解）]
由可控性结构分解，取坐标使
\begin{equation}
  \tilde{\mathbf A}=\begin{bmatrix}\tilde{\mathbf A}_{11}&\tilde{\mathbf A}_{12}\\\mathbf 0&\tilde{\mathbf A}_{22}\end{bmatrix},\qquad \tilde{\mathbf B}=\begin{bmatrix}\tilde{\mathbf B}_1\\\mathbf 0\end{bmatrix},
  \label{eq:syn-stab-decomp}
\end{equation}
其中 $(\tilde{\mathbf A}_{11},\tilde{\mathbf B}_1)$ 可控、$\tilde{\mathbf A}_{22}$ 是不可控子系统。引入 $\tilde{\mathbf K}=[\tilde{\mathbf K}_1\ \tilde{\mathbf K}_2]$，闭环阵
\begin{equation}
  \tilde{\mathbf A}-\tilde{\mathbf B}\tilde{\mathbf K}=\begin{bmatrix}\tilde{\mathbf A}_{11}-\tilde{\mathbf B}_1\tilde{\mathbf K}_1&\tilde{\mathbf A}_{12}-\tilde{\mathbf B}_1\tilde{\mathbf K}_2\\\mathbf 0&\tilde{\mathbf A}_{22}\end{bmatrix}.
  \label{eq:syn-stab-closed}
\end{equation}
块上三角，其特征值 $=\operatorname{eig}(\tilde{\mathbf A}_{11}-\tilde{\mathbf B}_1\tilde{\mathbf K}_1)\cup\operatorname{eig}(\tilde{\mathbf A}_{22})$。可控块 $(\tilde{\mathbf A}_{11},\tilde{\mathbf B}_1)$ 可经 $\tilde{\mathbf K}_1$ 任意配置（\cref{thm:syn-pole-placement}），故可全部配到左半平面；而 $\tilde{\mathbf A}_{22}$ 的特征值\emph{不受 $\tilde{\mathbf K}$ 影响}。于是整个系统可镇定 $\iff$ $\tilde{\mathbf A}_{22}$（不可控子系统）的特征值全具负实部，即不可控子系统渐近稳定。
\end{proof}

\begin{note}
\textbf{可镇定（stabilizable）的现代术语.}\;\cref{thm:syn-stabilize} 即“$(\mathbf A,\mathbf B)$ 可镇定”的定义——所有\emph{不稳定}模态可控（对偶于 \cref{ch:controllability-obsv} 的“可检测”=所有不稳定模态可观）。可镇定是 LQR / $H_\infty$ 解存在的最弱前提（比完全可控弱），导论 \cref{sec:ctrl-lqr} 的 CARE 存在唯一性正建立在“可镇定 + 可检测”之上；严格证明见 \cref{ch:lqr-lqg-riccati}。从若尔当型看：系统可镇定 $\iff$ 所有与\emph{不稳定}特征值相关的若尔当块都可控\cite{zhang2017modern}。
\end{note}

\paragraph{输出反馈镇定的局限。}
对偶地，从输出到 $\dot{\mathbf x}$ 反馈可镇定的充要条件是\emph{不可观子系统}渐近稳定（用观测器思路，见 \cref{sec:syn-observer}）。但\emph{静态输出反馈} $\mathbf u=-\mathbf H\mathbf y$ 的镇定能力弱得多：即便系统可控可观，也未必能用静态输出反馈镇定。一个经典反例\cite{liubao2006modern}：某可控可观系统经任意 $\mathbf H$ 反馈后特征多项式始终\emph{缺 $s^2$ 项}（某系数恒为零），故无论如何选 $\mathbf H$ 都配不出全左半平面极点。这正是“为何要状态反馈 + 观测器”而非满足于静态输出反馈的根本原因——后者的可达极点集只是一条低维曲面，而前者（经观测器）能覆盖全空间。

% ════════════════════════════════════════════════════════════════
\section{系统解耦}
\label{sec:syn-decoupling}

\paragraph{动机：把纠缠的多变量系统拆成独立单回路。}
多输入多输出系统的麻烦在于\emph{耦合}：动一个输入，好几个输出同时晃。\textbf{解耦（decoupling）}的目标是设计控制律，使每个输入只控制对应的\emph{一个}输出、每个输出只受对应的\emph{一个}输入支配——这样多变量系统就化为一组互不干扰的单变量子系统，可分别整定（如航空、化工的多回路控制）。形式化：寻找控制律使闭环传递函数阵 $\mathbf G(s;\mathbf K,\mathbf F)$ 成为\emph{非奇异对角}有理分式阵。

\paragraph{两条路线。}
\begin{itemize}
  \item \textbf{前馈补偿器解耦}：在待解耦系统前串接一个补偿器 $\mathbf W_d(s)$，使串联组合 $\mathbf W_0(s)\mathbf W_d(s)$ 成对角阵。只要 $\mathbf W_0(s)$ 满秩即可解耦，但\emph{增加系统维数}（补偿器是动态环节）。
  \item \textbf{状态反馈解耦}：用 $\mathbf u=\mathbf F\mathbf v-\mathbf K\mathbf x$（$\mathbf F$ 为 $m\times m$ 非奇异输入变换阵）。\emph{不增加维数}，但解耦条件苛刻得多。下面重点讲这条。
\end{itemize}

\paragraph{两个特征量。}
设 $\mathbf c_i$ 为输出阵 $\mathbf C$ 的第 $i$ 行。定义\cite{liubao2006modern,zhang2017modern}：
\begin{itemize}
  \item \textbf{解耦阶数} $d_i$：使 $\mathbf c_i\mathbf A^l\mathbf B\neq\mathbf 0$ 成立的\emph{最小}整数 $l\in\{0,1,\dots,n-1\}$\footnote{源（刘豹 §5.4.2 式 5.69）OCR 把此范围误印为 $0\le l\le m-1$（M2-b L2659），按 Falb--Wolovich 标准应为 $0\le l\le n-1$（$n$ 为系统维数），本书据勘误改正。};若对一切 $l$ 皆 $\mathbf c_i\mathbf A^l\mathbf B=\mathbf 0$，则 $d_i=n-1$。直观上 $d_i+1$ 是第 $i$ 个输出对输入的“相对阶”（输出需微分几次才显式出现输入）。
  \item \textbf{解耦矩阵} $\mathbf E$：$m\times m$ 阵，第 $i$ 行为 $\mathbf c_i\mathbf A^{d_i}\mathbf B$，即
  \begin{equation}
    \mathbf E=\begin{bmatrix}\mathbf c_1\mathbf A^{d_1}\mathbf B\\\mathbf c_2\mathbf A^{d_2}\mathbf B\\\vdots\\\mathbf c_m\mathbf A^{d_m}\mathbf B\end{bmatrix}.
    \label{eq:syn-E}
  \end{equation}
\end{itemize}

\begin{theorem}[可解耦判据]
\label{thm:syn-decouple}
受控系统 $\Sigma_0=(\mathbf A,\mathbf B,\mathbf C)$ 采用状态反馈 $\mathbf u=\mathbf F\mathbf v-\mathbf K\mathbf x$ 可实现解耦的\textbf{充要条件}是解耦矩阵 $\mathbf E$ 非奇异\cite{liubao2006modern,zhang2017modern}。此时取
\begin{equation}
  \mathbf F=\mathbf E^{-1},\qquad \mathbf K=\mathbf E^{-1}\mathbf F^*,\quad \mathbf F^*=\begin{bmatrix}\mathbf c_1\mathbf A^{d_1+1}\\\vdots\\\mathbf c_m\mathbf A^{d_m+1}\end{bmatrix},
  \label{eq:syn-decouple-KF}
\end{equation}
即得\textbf{积分型解耦系统}：每个子系统是一个 $(d_i+1)$ 阶积分器 $1/s^{d_i+1}$。
\end{theorem}

\begin{insight}
\textbf{为什么是 $\mathbf E$ 非奇异？}\;把第 $i$ 个输出微分 $d_i+1$ 次，第一次让输入显式出现，其系数恰是 $\mathbf c_i\mathbf A^{d_i}\mathbf B$（$\mathbf E$ 的第 $i$ 行）。要让“每个输出独立地被一个输入驱动”，需这 $m$ 行\emph{线性无关}——即 $\mathbf E$ 可逆，才能用 $\mathbf F=\mathbf E^{-1}$ 把交叉耦合\emph{精确抵消}。$\mathbf E$ 奇异，则某些输出的“首现输入方向”重合，无论怎么反馈都拆不开。
\end{insight}

\begin{example}[状态反馈解耦设计]
\label{ex:syn-decouple}
某双输入双输出系统经计算得 $d_1=d_2=1$，且解耦矩阵 $\mathbf E=\big[\begin{smallmatrix}1&0\\0&1\end{smallmatrix}\big]=\mathbf I$\cite{liubao2006modern}。因 $\mathbf E$ 非奇异，系统可状态反馈解耦。由 \cref{eq:syn-decouple-KF}：$\mathbf F=\mathbf E^{-1}=\mathbf I$，$\mathbf K=\mathbf E^{-1}\mathbf F^*=\mathbf F^*$（$\mathbf F^*$ 各行为 $\mathbf c_i\mathbf A^{d_i+1}=\mathbf c_i\mathbf A^2$）。代入控制律 $\mathbf u=\mathbf F\mathbf v-\mathbf K\mathbf x$ 后，闭环成两个独立的二阶积分器 $1/s^2$（$d_i+1=2$），即每个输入只驱动对应一个输出。$\mathbf K$ 各元素的物理作用，是\emph{抵消状态变量间的交叉耦合项}，从而实现自治控制。最后对这两个 $1/s^2$ 子系统分别追加状态反馈，把极点配到（如）$-1,-1$ 与 $-2,-2$，即得性能合格的解耦系统。
\end{example}

\paragraph{解耦后还需配极点。}
积分型解耦系统所有极点都在\emph{原点}（纯积分器），动态性能极差、且临界不稳。故解耦只是第一步：还须对每个已解耦的单变量子系统\emph{追加状态反馈}，把其极点配到期望位置（\cref{thm:syn-pole-placement}）。这一步要求各子系统可控——表面上解耦不涉及可控性，但“积分型解耦系统无实用价值、要靠极点配置救活”，故系统仍须可控（至少可镇定），否则解耦失去意义\cite{zhang2017modern}。完整设计流程：求 $d_i,\mathbf E$ $\to$ 验 $\mathbf E$ 非奇异 $\to$ 由 \cref{eq:syn-decouple-KF} 算 $(\mathbf K,\mathbf F)$ 得积分型 $\to$ 对各子系统追加极点配置。

\begin{note}
\textbf{能解耦标准形与前馈补偿器解耦（补全两条路线细节）.}\;
\emph{(i) 能解耦标准形.}\;积分型解耦系统经非奇异变换可化为\textbf{能解耦标准形} $\hat\Sigma=(\hat{\mathbf A},\hat{\mathbf B},\hat{\mathbf C})$：三阵均\emph{块对角}，第 $i$ 块是一条 $p_i=d_i+1$ 阶积分器链\cite{liubao2006modern}
\begin{equation}
  \hat{\mathbf A}=\operatorname{diag}(\mathbf A_1,\dots,\mathbf A_m),\ \ \mathbf A_i=\begin{bmatrix}\mathbf 0&\mathbf I_{d_i}\\0&\mathbf 0\end{bmatrix}_{p_i\times p_i};\quad
  \hat{\mathbf B}=\operatorname{diag}(\mathbf b_1,\dots,\mathbf b_m),\ \mathbf b_i=\begin{bmatrix}0\\\vdots\\0\\1\end{bmatrix};\quad
  \hat{\mathbf C}=\operatorname{diag}(\mathbf c_1,\dots,\mathbf c_m),\ \mathbf c_i=[1\ 0\cdots0].
  \label{eq:syn-decouple-canon}
\end{equation}
在此标准形下，“解耦 $+$ 任意配极点”的附加状态反馈 $(\hat{\mathbf K},\hat{\mathbf F})$ 也必\emph{块对角}：$\hat{\mathbf K}=\operatorname{diag}(\mathbf k_1,\dots,\mathbf k_m)$、$\hat{\mathbf F}=\operatorname{diag}(f_1,\dots,f_m)$（$f_i\neq0$），每块 $\mathbf k_i=[k_{i0}\ \cdots\ k_{id_i}]$ 独立地把第 $i$ 条积分器链的 $d_i+1$ 个极点配到期望位置、互不串扰。这把“解耦后逐子系统配极点”落成了\emph{块对角配置}，且 $\hat\Sigma$ 是闭环传函阵 $\mathbf W_{KF}(s)$ 的一个\emph{最小实现}。
\emph{(ii) 前馈补偿器解耦.}\;另一条路线（\cref{sec:syn-decoupling} 开头第一种）在对象前串接前馈补偿器 $\mathbf W_d(s)$，使串联组合 $\mathbf W(s)=\mathbf W_0(s)\mathbf W_d(s)$ 成对角阵 $\operatorname{diag}(W_{11},\dots,W_{mm})$。只要 $\mathbf W_0(s)$ 满秩（$\mathbf W_0^{-1}$ 存在），即可取
\begin{equation}
  \mathbf W_d(s)=\mathbf W_0^{-1}(s)\,\mathbf W(s),
  \label{eq:syn-decouple-ff}
\end{equation}
解耦后各子系统特性由对角元 $W_{ii}(s)$ 自由规定\cite{liubao2006modern}。代价是 $\mathbf W_d$ 为动态环节、\emph{增加系统维数}（对比状态反馈解耦不增维但条件苛刻）。\textbf{两条路线的统一}：状态反馈解耦其实是“串联补偿器退化为零阶常数阵”的特例；对不能用纯状态反馈解耦者，还可\emph{兼用}状态反馈 $+$ 串联补偿器去解耦更一般的对象。
\end{note}

\begin{pitfall}[解耦的代价：抵消耦合 = 对模型误差敏感]
状态反馈解耦靠 $\mathbf F=\mathbf E^{-1}$ \emph{精确抵消}交叉耦合——但这抵消依赖 $\mathbf E$（含 $\mathbf A,\mathbf B,\mathbf C$）的精确值。模型一旦有误差，抵消不彻底，残余耦合重新出现，且解耦控制器对此\emph{无鲁棒性补偿}。这是“完美解耦”与“鲁棒”之间的张力：越是把回路拆得干净，越是把设计绑死在精确模型上。工程上常\emph{不追求严格解耦}，而容忍弱耦合、换取鲁棒裕度——与 \cref{sec:syn-robust} 的鲁棒思想一脉相承。
\end{pitfall}

% ════════════════════════════════════════════════════════════════
\section{状态观测器}
\label{sec:syn-observer}

到此为止，极点配置、解耦、乃至下一章的最优控制，都建立在\emph{全状态反馈} $\mathbf u=-\mathbf K\mathbf x$ 之上。但有一道现实鸿沟：\textbf{状态变量并不都可直接测量}——有些因传感器昂贵或物理上不可达而测不到（如机械臂各关节的精确速度、电机内部磁链、足式机器人浮动基座的全速度）。于是问题尖锐起来：手里只有输入 $\mathbf u$ 与输出 $\mathbf y=\mathbf C\mathbf x$，如何重构出 $\mathbf x$，好把它喂给状态反馈？这就是\textbf{状态观测 / 状态重构}问题，由龙伯格（Luenberger）于 1964 年前后解决\cite{liubao2006modern,luenberger1971introduction}。

\paragraph{反面一：微分器重构——理论可行、工程不可行。}
最朴素的想法：既然 $\mathbf y=\mathbf C\mathbf x$，把 $\mathbf y$ 反复求导 $\dot{\mathbf y},\ddot{\mathbf y},\dots$ 凑够 $n$ 个方程，若 $(\mathbf A,\mathbf C)$ 可观，可观性矩阵 $\mathcal O=[\mathbf C;\mathbf C\mathbf A;\cdots;\mathbf C\mathbf A^{n-1}]$ 满秩，便能从 $\{\mathbf y,\dot{\mathbf y},\dots,\mathbf y^{(n-1)}\}$ 解出 $\mathbf x$。理论上无懈可击——\emph{可观即可由输出及其各阶导数估出状态}。但这构造\textbf{毫无工程价值}\cite{liubao2006modern}：它需要 $0$ 到 $n-1$ 阶微分器，而\emph{微分会剧烈放大测量噪声}（微分的频率响应 $\propto j\omega$，高频噪声被成倍灌大）。实测信号总带噪，逐阶微分后状态估值会被噪声淹没。\textbf{这是整条观测器理论的出发点：必须找一个不依赖微分的重构办法。}

\paragraph{反面二：开环观测器——只是模型副本，会漂。}
第二个想法：既然知道模型 $(\mathbf A,\mathbf B)$，就在线跑一个\emph{副本} $\dot{\hat{\mathbf x}}=\mathbf A\hat{\mathbf x}+\mathbf B\mathbf u$，用它的 $\hat{\mathbf x}$ 当估计。这叫\emph{开环观测器}。误差 $\tilde{\mathbf x}=\mathbf x-\hat{\mathbf x}$ 满足 $\dot{\tilde{\mathbf x}}=\mathbf A\tilde{\mathbf x}$——若 $\mathbf A$ 稳定则误差自行衰减，但\emph{衰减快慢由 $\mathbf A$ 锁死、不可调}；若 $\mathbf A$ 不稳则误差发散。更糟的是它\emph{没用上输出信息 $\mathbf y$}：初值不符、扰动、参数失配都会让副本与真实状态越漂越远，且无从纠正。开环观测器没有实用意义。

\paragraph{核心思想：模型副本 + 输出残差纠正 = 渐近观测器。}
把两个反面合起来想：副本提供“预测”，输出提供“纠错”。当 $\hat{\mathbf x}\neq\mathbf x$ 时，估计输出 $\hat{\mathbf y}=\mathbf C\hat{\mathbf x}$ 也偏离实测 $\mathbf y$，产生\emph{残差} $\mathbf y-\mathbf C\hat{\mathbf x}$；把这残差经增益 $\mathbf L$ 反馈进副本的每个积分器，去推动 $\hat{\mathbf x}$ 趋近 $\mathbf x$。这便是\textbf{龙伯格（渐近）观测器}（导论 \cref{eq:ctrl-luenberger}）：
\begin{equation}
  \dot{\hat{\mathbf x}}=\underbrace{\mathbf A\hat{\mathbf x}+\mathbf B\mathbf u}_{\text{模型预测}}+\underbrace{\mathbf L(\mathbf y-\mathbf C\hat{\mathbf x})}_{\text{输出残差纠正}}.
  \label{eq:syn-luenberger}
\end{equation}
它只用 $\mathbf u,\mathbf y$ 与代数运算（无微分），克服了反面一；用残差闭环纠错（用上 $\mathbf y$），克服了反面二。

\begin{figure}[H]
\centering
\begin{tikzpicture}[>=Stealth, font=\small, node distance=7mm,
  box/.style={draw, minimum height=7mm, minimum width=18mm, fill=second!8},
  sum/.style={draw, circle, inner sep=1pt, minimum size=4mm}]
  % plant
  \node[box, fill=main!8] (plant) at (0,0) {$\Sigma_0:\ \dot{\mathbf x}=\mathbf A\mathbf x+\mathbf B\mathbf u$};
  \node[right=18mm of plant] (y) {$\mathbf y$};
  \draw[->] (plant)--(y) node[midway,above]{$\mathbf C\mathbf x$};
  % observer
  \node[box, below=16mm of plant] (obs) {$\dot{\hat{\mathbf x}}=\mathbf A\hat{\mathbf x}+\mathbf B\mathbf u+\mathbf L\mathbf r$};
  \node[sum, right=14mm of obs] (s) {$+$};
  \node[box, below=8mm of s, minimum width=10mm] (Cobs) {$\mathbf C$};
  \node[left=10mm of plant] (u) {$\mathbf u$};
  \draw[->] (u)--(plant);
  \draw[->] (u) |- (obs);
  % residual
  \coordinate (ytap) at ($(y)+(0,0)$);
  \draw[->] (ytap) -- ++(0,-8mm) -| (s) node[pos=0.95,right]{$\mathbf y$};
  \draw[->] (obs.east) -- (s) node[midway,above]{$\hat{\mathbf x}$};
  \draw[->] (obs.south) |- (Cobs);
  \draw[->] (Cobs) -| (s) node[pos=0.9,left]{$-\hat{\mathbf y}$};
  \draw[->] (s) -- ++(0,8mm) -| node[pos=0.92,right]{$\mathbf r=\mathbf y-\mathbf C\hat{\mathbf x}$} ($(obs.north)+(0,0)$);
\end{tikzpicture}
\caption{龙伯格渐近观测器：模型副本（下框）以 $\mathbf u,\mathbf y$ 为输入，用输出残差 $\mathbf r=\mathbf y-\mathbf C\hat{\mathbf x}$ 经增益 $\mathbf L$ 纠正预测，使 $\hat{\mathbf x}$ 渐近逼近真实状态 $\mathbf x$。（仿刘豹 §5.5 图 5.16 重绘）}
\label{fig:syn-observer}
\end{figure}

\paragraph{误差动态：观测器设计的核心方程。}
用 \cref{eq:syn-luenberger} 减真实动力学 $\dot{\mathbf x}=\mathbf A\mathbf x+\mathbf B\mathbf u$（注意 $\mathbf B\mathbf u$ 项相消、输入消失）：
\begin{align}
  \dot{\tilde{\mathbf x}}&=\dot{\mathbf x}-\dot{\hat{\mathbf x}}=\mathbf A\mathbf x-\big[\mathbf A\hat{\mathbf x}+\mathbf L(\mathbf C\mathbf x-\mathbf C\hat{\mathbf x})\big]\notag\\
  &=\mathbf A(\mathbf x-\hat{\mathbf x})-\mathbf L\mathbf C(\mathbf x-\hat{\mathbf x})=(\mathbf A-\mathbf L\mathbf C)\tilde{\mathbf x}.
  \label{eq:syn-obs-error}
\end{align}
\textbf{误差动态由 $(\mathbf A-\mathbf L\mathbf C)$ 唯一决定、不依赖输入与状态}（与导论 \cref{eq:ctrl-obs-error} 同）。于是：若能选 $\mathbf L$ 使 $(\mathbf A-\mathbf L\mathbf C)$ 的特征值全具负实部，则 $\tilde{\mathbf x}\to\mathbf 0$，$\hat{\mathbf x}$ 渐近收敛于 $\mathbf x$；衰减速度由这些特征值（误差极点）决定。若初值恰好 $\tilde{\mathbf x}(0)=\mathbf 0$，则 $\tilde{\mathbf x}\equiv\mathbf 0$、估计严格等于真值。

\begin{theorem}[观测器增益设计 = 对偶极点配置]
\label{thm:syn-obs-design}
若 $(\mathbf A,\mathbf C)$ 完全可观，则可选观测器增益 $\mathbf L$ 把 $(\mathbf A-\mathbf L\mathbf C)$ 的特征值\emph{任意}配置。设计经\textbf{对偶}转化为一个极点配置问题：$(\mathbf A-\mathbf L\mathbf C)$ 与其转置 $(\mathbf A^\top-\mathbf C^\top\mathbf L^\top)$ 同特征值，而由对偶原理（\cref{thm:co-duality}），$(\mathbf A,\mathbf C)$ 可观 $\iff$ $(\mathbf A^\top,\mathbf C^\top)$ 可控\cite{liubao2006modern}。故对“对偶系统” $(\mathbf A^\top,\mathbf C^\top)$ 设计状态反馈 $\mathbf L^\top$（用 \cref{thm:syn-pole-placement} / Ackermann），转置回来即得 $\mathbf L$。
\end{theorem}

\begin{note}
\textbf{对偶转置中不要凭空多出导数点.}\;源（刘豹 §5.2.3 定理证明）在写对偶变换 $(\mathbf A^\top+\mathbf c^\top\mathbf G^\top)^\top=\mathbf A+\mathbf G\mathbf c$ 时，OCR 误把 $\mathbf G$ 识成 $\dot{\mathbf G}$（M2-b L2322）——这里\emph{没有任何导数}，纯是矩阵转置 $(\mathbf c^\top\mathbf G^\top)^\top=\mathbf G\mathbf c$。本书据勘误写作无导数点形式。
\end{note}

\begin{theorem}[观测器存在的充要条件]
\label{thm:syn-obs-exist}
对 $\Sigma_0=(\mathbf A,\mathbf B,\mathbf C)$，状态观测器存在的充要条件是其\textbf{不可观子系统渐近稳定}\cite{liubao2006modern,zhang2017modern}。
\end{theorem}

\noindent 这是 \cref{thm:syn-stabilize}（镇定）的\emph{对偶}：镇定要求“不可控子系统稳”，观测器存在要求“不可观子系统稳”。若系统完全可观，则不可观子系统为空、条件自动满足，且 $\mathbf L$ 可任意配；若仅不完全可观但不可观子系统稳，仍可造观测器，但 $\hat{\mathbf x}\to\mathbf x$ 的速度受不可观子系统极点限制、不可任意加快。“完全可观 $\Rightarrow$ 观测器可任意配极点”正是“完全可控 $\Rightarrow$ 极点可任意配”的镜像。

\begin{proof}[证明要点（对偶到 \cref{thm:syn-stabilize}）]
把可观性结构分解（\cref{sec:co-decomp} 的对偶）施于 $(\mathbf A,\mathbf C)$，取坐标使可观块与不可观块分离，不可观子系统记 $\mathbf A_{\bar o}$。观测器误差阵 $(\mathbf A-\mathbf L\mathbf C)$ 在该坐标下亦呈块三角，其谱中\emph{不可观块对应的特征值不受 $\mathbf L$ 影响}（残差 $\mathbf y-\mathbf C\hat{\mathbf x}$ 看不到不可观模态，纠不了它）。故误差可全部配到左半平面 $\iff$ $\mathbf A_{\bar o}$ 本身已稳定。这与 \cref{thm:syn-stabilize} 的证明逐字对偶（把 $\mathbf B,\mathbf K,\tilde{\mathbf A}_{22}$ 换成 $\mathbf C^\top,\mathbf L^\top,\mathbf A_{\bar o}$）。
\end{proof}

\begin{insight}
\textbf{观测器极点“配得快”的隐性代价。}\;\cref{eq:syn-obs-error} 说误差衰减率由 $(\mathbf A-\mathbf L\mathbf C)$ 的极点决定，似乎配得越左、收敛越快越好。但快收敛要求\emph{大} $\mathbf L$：残差被高增益放大灌进 $\hat{\mathbf x}$，测量噪声也被同等放大——估计虽“快”却“抖”。这与极点配置的“快但费力”是对偶的两面：那边是控制量代价，这边是噪声代价。\textbf{卡尔曼滤波的全部意义，正是用噪声协方差把这个“快 vs 抖”的折中\emph{最优化}}（\cref{ch:eskf}），而非靠手配极点拍脑袋。这就是为什么确定性观测器（本章）讲“会设计”，而最优观测器（卡尔曼）要单列一章讲“最优折中”。
\end{insight}

\begin{pitfall}[把“可观”当成观测器存在的必要条件]
常见误解：“要造观测器，系统必须完全可观。”实则\emph{可观是充分、不是必要}——\cref{thm:syn-obs-exist} 的充要条件是更弱的“不可观子系统稳定”（即\emph{可检测}）。一个不完全可观的系统，只要那些看不见的模态本就自行衰减，仍可造出收敛的观测器（只是收敛速度受这些模态极点限制、不能任意加快）。这与“镇定只需可镇定、不必完全可控”完全对偶。把“可观”误当必要条件，会让你对一大类\emph{可检测但不可观}的实际系统（带冗余传感器、或某些内部状态恒稳的机电系统）错误地放弃观测器设计。
\end{pitfall}

\begin{example}[全维观测器设计]
\label{ex:syn-observer}
设 $W_0(s)=\dfrac{1}{s(s+6)}$，已知系统可控可观。设计全维观测器使误差极点为 $-10,-10$\cite{liubao2006modern}。

\textbf{解.}\;写可观标准 II 型实现 $\dot{\mathbf x}=\big[\begin{smallmatrix}0&0\\1&-6\end{smallmatrix}\big]\mathbf x+\big[\begin{smallmatrix}1\\0\end{smallmatrix}\big]u$\footnote{源（刘豹例 5-11）OCR 把 $\mathbf b$ 误印为 $[0;1]$；多重交叉验证（$\mathbf A+\mathbf b\mathbf K$ 结构、特征式、观测器方程）确认应为 $\mathbf b=[1;0]$（M2-c L3747），本书据勘误改正。},$\mathbf c=[0\ 1]$。令 $\mathbf L=[\ell_1\ \ell_2]^\top$，则 $(\mathbf A-\mathbf L\mathbf c)$ 的特征多项式应等于 $(s+10)^2=s^2+20s+100$，比对系数解出 $\ell_1,\ell_2$。代入 \cref{eq:syn-luenberger} 即得观测器方程 $\dot{\hat{\mathbf x}}=\mathbf A\hat{\mathbf x}+\mathbf b u+\mathbf L(y-\mathbf c\hat{\mathbf x})$。
\end{example}

% ════════════════════════════════════════════════════════════════
\section{降维观测器}
\label{sec:syn-reduced}

\paragraph{动机：测得到的状态何必再估？}
全维观测器重构\emph{全部} $n$ 个状态——可输出 $\mathbf y=\mathbf C\mathbf x$ 本身已提供了 $m$ 个状态的（线性组合）信息，何必把这 $m$ 维也重新估一遍？若 $\operatorname{rank}\mathbf C=m$，则 $\mathbf y$ 直接给出 $m$ 个状态分量，只需对其余 $(n-m)$ 个\emph{未测}分量造一个 $(n-m)$ 维观测器即可。这就是\textbf{降维（Luenberger 降阶）观测器}，维数从 $n$ 降到 $n-m$，省算且省噪声通道\cite{liubao2006modern,zhang2017modern}。

\paragraph{两步构造。}
\emph{第一步：按可测性分解状态。}\;经线性变换 $\bar{\mathbf x}=\mathbf P\mathbf x$ 把状态分成两块
\begin{equation}
  \bar{\mathbf x}=\begin{bmatrix}\bar{\mathbf x}_1\\\bar{\mathbf x}_2\end{bmatrix},\qquad \bar{\mathbf x}_2=\bar{\mathbf y}\ (\text{可由输出直接测得}),\quad \bar{\mathbf x}_1\in\mathbb R^{\,n-m}\ (\text{待重构}),
  \label{eq:syn-red-split}
\end{equation}
变换后系统写成分块形式
\begin{equation}
  \begin{bmatrix}\dot{\bar{\mathbf x}}_1\\\dot{\bar{\mathbf x}}_2\end{bmatrix}=\begin{bmatrix}\bar{\mathbf A}_{11}&\bar{\mathbf A}_{12}\\\bar{\mathbf A}_{21}&\bar{\mathbf A}_{22}\end{bmatrix}\begin{bmatrix}\bar{\mathbf x}_1\\\bar{\mathbf x}_2\end{bmatrix}+\begin{bmatrix}\bar{\mathbf B}_1\\\bar{\mathbf B}_2\end{bmatrix}\mathbf u.
  \label{eq:syn-red-block}
\end{equation}

\emph{第二步：对 $\bar{\mathbf x}_1$ 造 $(n-m)$ 维观测器。}\;把 $\bar{\mathbf x}_1$ 视作子系统状态，$\bar{\mathbf x}_2=\bar{\mathbf y}$ 已知，则 $\bar{\mathbf x}_1$ 的动态为 $\dot{\bar{\mathbf x}}_1=\bar{\mathbf A}_{11}\bar{\mathbf x}_1+(\bar{\mathbf A}_{12}\bar{\mathbf y}+\bar{\mathbf B}_1\mathbf u)$，其中括号项\emph{已知}（充当“输入”）；而 $\dot{\bar{\mathbf y}}-\bar{\mathbf A}_{22}\bar{\mathbf y}-\bar{\mathbf B}_2\mathbf u=\bar{\mathbf A}_{21}\bar{\mathbf x}_1$ 充当“输出方程”，$\bar{\mathbf A}_{21}$ 是子系统输出阵。由 $(\mathbf A,\mathbf C)$ 可观可推 $(\bar{\mathbf A}_{11},\bar{\mathbf A}_{21})$ 也可观，故 $\bar{\Sigma}_1$ 存在观测器。仿全维做法得观测器方程；为\emph{避免对 $\bar{\mathbf y}$ 求导}（又会引入微分放噪），引入辅助变量 $\mathbf w=\hat{\bar{\mathbf x}}_1-\bar{\mathbf L}\bar{\mathbf y}$，把含 $\dot{\bar{\mathbf y}}$ 的项吸收掉，最终估计为
\begin{equation}
  \hat{\mathbf x}=\begin{bmatrix}\hat{\bar{\mathbf x}}_1\\\bar{\mathbf x}_2\end{bmatrix}=\begin{bmatrix}\mathbf w+\bar{\mathbf L}\bar{\mathbf y}\\\bar{\mathbf y}\end{bmatrix}=\begin{bmatrix}\mathbf I\\\mathbf 0\end{bmatrix}\mathbf w+\begin{bmatrix}\bar{\mathbf L}\\\mathbf I\end{bmatrix}\bar{\mathbf y}.
  \label{eq:syn-red-est}
\end{equation}
\footnote{源（刘豹式 5.110）OCR 在此处多出三处负号（把上划线 $\bar{\mathbf x}_2,\bar{\mathbf y}$ 误识为前导负号，并把底块 $\mathbf I$ 误成 $-\mathbf I$，M2-c L3332）；按 $\bar{\mathbf x}_2=\bar{\mathbf y}$ 可测、底块系数取 $+\mathbf I$，本书据勘误改正。}
观测器内部 $\dot{\mathbf w}=(\bar{\mathbf A}_{11}-\bar{\mathbf L}\bar{\mathbf A}_{21})\mathbf w+(\cdots)$，误差极点由 $(\bar{\mathbf A}_{11}-\bar{\mathbf L}\bar{\mathbf A}_{21})$ 配置，仍是 $(n-m)$ 维的对偶极点配置。

\begin{insight}
\textbf{降维 vs 全维的取舍。}\;降维观测器维数低、计算省，且\emph{直接用上可测分量}（不经重构、无延迟、无估计误差）。但它对\emph{测量噪声}更敏感：$\mathbf y$ 中的噪声直接进入 $\hat{\mathbf x}$ 的 $\bar{\mathbf x}_2$ 块、不经任何滤波，而全维观测器把 $\mathbf y$ 也当“带噪测量”过一遍模型平滑。故在噪声大的场合，全维（乃至卡尔曼滤波，\cref{ch:eskf}）反而更稳。这是“维数 $\leftrightarrow$ 抗噪”的工程权衡，没有放之四海皆准的赢家。
\end{insight}

% ════════════════════════════════════════════════════════════════
\section{利用观测器实现状态反馈：分离原理}
\label{sec:syn-separation}

\paragraph{动机：把估计的状态喂回反馈，闭环还稳吗？}
观测器解决了状态重构，于是可把状态反馈中的真状态换成估计：$\mathbf u=-\mathbf K\hat{\mathbf x}$。这便是\textbf{基于观测器的输出反馈}——“先估计、再控制”。但一个尖锐的疑问随之而来：原本基于真状态 $\mathbf x$ 设计的 $\mathbf K$（保证 $\mathbf A-\mathbf B\mathbf K$ 稳定），现在喂的是\emph{带误差}的 $\hat{\mathbf x}=\mathbf x-\tilde{\mathbf x}$，控制器极点与观测器极点会不会耦合、互相破坏？答案是一个出奇优美的结果。

\begin{figure}[H]
\centering
\begin{tikzpicture}[>=Stealth, font=\small,
  box/.style={draw, minimum height=7mm, minimum width=20mm, fill=main!7},
  obsbox/.style={draw, minimum height=7mm, minimum width=20mm, fill=second!8},
  sum/.style={draw, circle, inner sep=1pt, minimum size=4mm}]
  \node[sum] (s1) at (0,0) {$+$};
  \node[box, right=10mm of s1] (plant) {$\Sigma_0=(\mathbf A,\mathbf B,\mathbf C)$};
  \node[right=14mm of plant] (y) {$\mathbf y$};
  \node[obsbox, below=11mm of plant] (obs) {观测器 $\hat{\mathbf x}$};
  \node[box, left=10mm of obs, minimum width=8mm] (K) {$-\mathbf K$};
  \node[left=8mm of s1] (v) {$\mathbf v$};
  \draw[->] (v)--(s1);
  \draw[->] (s1)--(plant) node[midway,above]{$\mathbf u$};
  \draw[->] (plant)--(y);
  \draw[->] (y) -- ++(0,-11mm) -| (obs.east) node[pos=0.05,right]{$\mathbf y$};
  \draw[->] (plant.south) -- ++(0,-4mm) -| (obs.north) node[pos=0.05,left]{$\mathbf u$};
  \draw[->] (obs)--(K) node[midway,above]{$\hat{\mathbf x}$};
  \draw[->] (K) |- (s1) node[pos=0.9,left]{};
\end{tikzpicture}
\caption{带全维观测器的状态反馈系统：观测器以 $\mathbf u,\mathbf y$ 重构 $\hat{\mathbf x}$，反馈律 $\mathbf u=\mathbf v-\mathbf K\hat{\mathbf x}$。分离原理保证控制器与观测器可独立设计。（仿刘豹 §5.6 图 5.21 重绘）}
\label{fig:syn-separation}
\end{figure}

\begin{theorem}[分离原理（闭环极点的分离性）]
\label{thm:syn-separation}
对 LTI 系统，用全维观测器估计的状态做反馈 $\mathbf u=\mathbf v-\mathbf K\hat{\mathbf x}$，则以 $(\mathbf x,\tilde{\mathbf x})$ 为状态的闭环系统矩阵\emph{块上三角}，其特征值是\textbf{控制器极点} $\operatorname{eig}(\mathbf A-\mathbf B\mathbf K)$ 与\textbf{观测器极点} $\operatorname{eig}(\mathbf A-\mathbf L\mathbf C)$ 之\emph{并}。因此只要 $(\mathbf A,\mathbf B,\mathbf C)$ 可控且可观，$\mathbf K$ 与 $\mathbf L$ 可\textbf{完全独立}设计\cite{liubao2006modern,astrom2008feedback}。
\end{theorem}

\noindent{}（这是导论 \cref{prop:ctrl-separation} 的展开版；导论给结论与分块证明，本章把它放进“先估计、再控制”的完整动机链中。）

\begin{proof}[证明（坐标换 $(\mathbf x,\tilde{\mathbf x})$ 致块上三角）]
取 $\mathbf u=\mathbf v-\mathbf K\hat{\mathbf x}=\mathbf v-\mathbf K(\mathbf x-\tilde{\mathbf x})$，代入 $\dot{\mathbf x}=\mathbf A\mathbf x+\mathbf B\mathbf u$：
\begin{equation}
  \dot{\mathbf x}=(\mathbf A-\mathbf B\mathbf K)\mathbf x+\mathbf B\mathbf K\tilde{\mathbf x}+\mathbf B\mathbf v;
\end{equation}
误差动态由 \cref{eq:syn-obs-error} 为 $\dot{\tilde{\mathbf x}}=(\mathbf A-\mathbf L\mathbf C)\tilde{\mathbf x}$（不含 $\mathbf v$）。合并成增广系统：
\begin{equation}
  \begin{bmatrix}\dot{\mathbf x}\\\dot{\tilde{\mathbf x}}\end{bmatrix}=\begin{bmatrix}\mathbf A-\mathbf B\mathbf K&\mathbf B\mathbf K\\\mathbf 0&\mathbf A-\mathbf L\mathbf C\end{bmatrix}\begin{bmatrix}\mathbf x\\\tilde{\mathbf x}\end{bmatrix}+\begin{bmatrix}\mathbf B\\\mathbf 0\end{bmatrix}\mathbf v.
  \label{eq:syn-sep-augmented}
\end{equation}
块上三角阵的特征值等于两对角块特征值之并，故闭环极点 $=\operatorname{eig}(\mathbf A-\mathbf B\mathbf K)\cup\operatorname{eig}(\mathbf A-\mathbf L\mathbf C)$。\emph{两组极点互不影响}，对应的特征多项式恰是两块特征多项式之积。证毕。
\end{proof}

\begin{insight}
\textbf{分离原理是“先估计、再控制”的理论许可证。}\;它宣告：状态估计器（观测器 / 卡尔曼，\cref{ch:eskf}）与状态反馈控制器（本章 / LQR，\cref{ch:optimal-control-basics}）\emph{可以分开设计}，拼起来仍稳定——这把一个 $2n$ 维耦合设计问题，干净地拆成两个 $n$ 维独立子问题。把“最优观测器（卡尔曼 / LQE）”与“最优状态反馈（LQR）”按此原理组合，即得 \textbf{LQG（线性二次高斯）}最优输出反馈控制器（导论 \cref{sec:ctrl-lqr} 的对偶半边）。\textbf{边界（重要）}：分离原理的“极点可分别配”对 LTI \emph{严格}成立，但 LQG 的\emph{鲁棒裕度不可分离}（Doyle 反例：LQR 有 $60^\circ$ 相位裕度，串上卡尔曼后裕度可降至任意小）——见 \cref{sec:syn-robust} 与 \cref{ch:robust-hinf}；非线性 / 受约束系统也只在局部近似成立。
\end{insight}

\paragraph{两个推论：传递函数不变性与等效性。}
带观测器状态反馈系统还有两条性质\cite{liubao2006modern}：
\begin{itemize}
  \item \textbf{传递函数不变性}：闭环传递函数阵\emph{与是否采用观测器无关}，等于直接状态反馈系统的传递函数。从 \cref{eq:syn-sep-augmented} 看：输入 $\mathbf v$ 经 $\big[\begin{smallmatrix}\mathbf B\\\mathbf 0\end{smallmatrix}\big]$ 进入，\emph{完全不激发}误差子空间 $\tilde{\mathbf x}$（其输入通道为 $\mathbf 0$）；故从 $\mathbf v$ 到 $\mathbf y$ 的传递只“看见” $(\mathbf A-\mathbf B\mathbf K,\mathbf B,\mathbf C)$ 这一块。等价地说，观测器极点 $\operatorname{eig}(\mathbf A-\mathbf L\mathbf C)$ 在传递函数中被\emph{零点精确对消}（它们对应不可控的误差子空间 $\tilde{\mathbf x}$）。因此该闭环虽\emph{不完全可控}（误差子系统不可控），但不可控部分恰是误差 $\tilde{\mathbf x}\to\mathbf 0$，不影响正常工作。
  \item \textbf{稳态等效性}：带观测器系统仅当 $t\to\infty$（$\tilde{\mathbf x}\to\mathbf 0$）才与直接状态反馈\emph{完全等价}；暂态期间因 $\hat{\mathbf x}\neq\mathbf x$ 而有差异（误差以 $\mathrm e^{(\mathbf A-\mathbf L\mathbf C)t}$ 衰减，经 $\mathbf B\mathbf K\tilde{\mathbf x}$ 项扰动 $\mathbf x$ 的动态）。可通过把观测器极点配得\emph{更快}（经验上比控制器极点快 $3\sim5$ 倍）来压缩暂态、加速 $\hat{\mathbf x}\to\mathbf x$——这是工程整定的常用法则（“观测器要比控制器快”），但配得过快会放大噪声（\cref{sec:syn-observer} 陷阱），故“快多少”本身是个折中。
\end{itemize}

\begin{insight}
\textbf{从一个 $2n$ 维难题到两个 $n$ 维易题。}\;\cref{eq:syn-sep-augmented} 的块上三角\emph{不是巧合}，而是“坐标选对了”的结果——选 $(\mathbf x,\tilde{\mathbf x})$ 而非 $(\mathbf x,\hat{\mathbf x})$，误差动态自我封闭（不被 $\mathbf x$ 反向耦合）。这把一个看似 $2n$ 维耦合的输出反馈设计，\emph{解耦}成“配 $\mathbf A-\mathbf B\mathbf K$（控制器）”与“配 $\mathbf A-\mathbf L\mathbf C$（观测器）”两个独立的 $n$ 维子问题。分离原理之所以深刻，正在于它揭示：\textbf{估计与控制在 LTI 下是\emph{可分离}的两件事}——这一论断在非线性 / 随机 / 受约束情形是否仍成立，构成了 LQG、MPC、乃至现代学习控制中反复出现的核心问题（\cref{ch:lqr-lqg-riccati}）。
\end{insight}

\begin{example}[综合设计：配极点 + 全维 / 降维观测器]
\label{ex:syn-full}
$W_0(s)=\dfrac{1}{s(s+6)}$，用状态反馈把闭环极点配到 $-4\pm j6$，并设计实现该反馈的全维与降维观测器（观测器极点 $-10,-10$）\cite{liubao2006modern}。

\textbf{解.}\;系统可控可观，存在反馈与观测器，据分离原理\emph{分别}设计：(i) 由 \cref{thm:syn-pole-placement} 求 $\mathbf K$ 使 $\operatorname{eig}(\mathbf A-\mathbf b\mathbf K)=\{-4\pm j6\}$；(ii) 由 \cref{thm:syn-obs-design} 求全维 $\mathbf L$ 使 $\operatorname{eig}(\mathbf A-\mathbf L\mathbf c)=\{-10,-10\}$；(iii) 由 \cref{sec:syn-reduced} 求 $1$ 维降维观测器（$n-m=2-1=1$）。最终控制律 $\mathbf u=\mathbf v-\mathbf K\hat{\mathbf x}$，三者拼装即成。这个例子完整走通了本章主线：可控 $\to$ 配极点，可观 $\to$ 造观测器，分离原理 $\to$ 拼起来稳。
\end{example}

\subsection{观测器即动态补偿器：与带补偿器输出反馈的等价性}
\label{sec:syn-compensator}

\paragraph{动机：观测器把“静态输出反馈”升级成了“动态输出反馈”。}
回到 \cref{sec:syn-output-pole,sec:syn-stabilize} 留下的悬念：\emph{静态}输出反馈 $\mathbf u=\mathbf v-\mathbf H\mathbf y$（无记忆、增益为常阵）一般\emph{配不动}全部 $n$ 个极点，甚至存在可控可观却无法被任何静态 $\mathbf H$ 镇定的系统。可基于观测器的反馈 $\mathbf u=\mathbf v-\mathbf K\hat{\mathbf x}$ 同样\emph{只用可测的} $\mathbf y$（它也是一种输出反馈！），却能恢复完整的极点配置能力。差别究竟在哪？答案是：观测器反馈\emph{不再是无记忆的}——它携带一个内部状态 $\hat{\mathbf x}$，是一种\textbf{动态输出反馈}。本节把这一点精确化：把“观测器 $+$ 状态反馈”当作从 $\mathbf y$ 到 $\mathbf u$ 的一个整体，它是一个 $n$ 阶\emph{动态补偿器}；并（循刘豹 §5.6.3）证明它在传递特性意义下\emph{完全等效}于一个经典的“串联补偿器 $+$ 反馈补偿器”输出反馈系统。

\begin{definition}[动态补偿器]
\label{def:syn-dyn-comp}
\cref{sec:syn-feedback} 的三种基本反馈（状态 / 输出 / 到 $\dot{\mathbf x}$）有个共同点：\emph{不引入新状态变量}，开环与闭环\emph{同维}，且反馈增益为\emph{常阵}。当需要更强的整形能力时，可在回路中引入一个\emph{带自身状态}的动态子系统，称为\textbf{动态补偿器（dynamic compensator）}\cite{liubao2006modern}。它与受控对象有两种接法：\emph{串联}连接（置于前向通道）与\emph{反馈}连接（置于反馈通道）。闭环维数 $=$ 受控对象维数 $+$ 补偿器维数。\textbf{最典型的动态补偿器，正是带状态观测器的状态反馈系统}——观测器即那个“带自身状态 $\hat{\mathbf x}$”的动态子系统。
\end{definition}

\begin{theorem}[动态补偿器可突破静态输出反馈的极点配置极限]
\label{thm:syn-dyn-comp-pole}
对完全可控的单输入单输出系统 $\Sigma_0=(\mathbf A,\mathbf b,\mathbf c)$，通过\emph{带动态补偿器的输出反馈}实现极点\emph{任意}配置的充要条件是：(1) $\Sigma_0$ 完全可观；(2) 动态补偿器的阶数为 $n-1$\cite{liubao2006modern}。若并不要求“任意”配置（只配部分极点），补偿器阶数可进一步降低。
\end{theorem}

\begin{insight}
\textbf{动态补偿器是经典“校正网络”的现代对应。}\;静态输出反馈只能让闭环极点沿\emph{根轨迹} $h\mathbf W_0(s)=-1$ 滑动（\cref{sec:syn-stabilize} 的反例正源于此），落不到任意指定位置。经典控制对此的补救，是\emph{串入校正网络（超前 / 滞后）}，靠\emph{增添开环零极点}改变根轨迹走向，使其穿过期望点\cite{liubao2006modern}。\cref{thm:syn-dyn-comp-pole} 是这一思想的状态空间版：一个 $n-1$ 阶动态补偿器恰好提供了“任意配 $n$ 个极点”所需的自由度。而\textbf{观测器 $+$ 状态反馈给出的，正是一个 $n$ 阶}（全维）\textbf{这样的动态补偿器}——比最小阶 $n-1$ 多一阶，这多出的一阶是“模块化分离设计”（\cref{thm:syn-separation}）的代价；用降维观测器（\cref{sec:syn-reduced}）则可把补偿器降到 $n-m$，逼近乃至达到最小阶（见本节末例）。（其闭环零点亦有定论：串联接法 $=$ 对象零点 $+$ 补偿器零点，反馈接法 $=$ 对象零点 $+$ 补偿器极点\cite{liubao2006modern}。）
\end{insight}

\paragraph{观测器 $+$ 状态反馈 $=$ 一个 $n$ 阶动态输出反馈补偿器。}
把龙伯格观测器（\cref{eq:syn-luenberger}）与反馈律 $\mathbf u=\mathbf v-\mathbf K\hat{\mathbf x}$ 合并，消去 $\mathbf u$：
\begin{equation}
  \dot{\hat{\mathbf x}}=(\mathbf A-\mathbf L\mathbf C)\hat{\mathbf x}+\mathbf L\mathbf y+\mathbf B(\mathbf v-\mathbf K\hat{\mathbf x})=(\mathbf A-\mathbf L\mathbf C-\mathbf B\mathbf K)\hat{\mathbf x}+\mathbf L\mathbf y+\mathbf B\mathbf v,\qquad \mathbf u=-\mathbf K\hat{\mathbf x}+\mathbf v.
  \label{eq:syn-comp-state}
\end{equation}
这是一个\emph{自包含}的动力系统：以测量 $\mathbf y$ 与参考 $\mathbf v$ 为输入、以 $\hat{\mathbf x}\in\mathbb R^n$ 为内部状态、以控制量 $\mathbf u$ 为输出。它\textbf{就是一个动态输出反馈补偿器}——给它测量，它\emph{动态地}算出控制。令 $\mathbf v=\mathbf 0$，得从 $\mathbf y$ 到 $\mathbf u$ 的传递阵：
\begin{equation}
  \mathbf u=-\mathbf C_K(s)\,\mathbf y,\qquad \mathbf C_K(s)=\mathbf K\,[\,s\mathbf I-(\mathbf A-\mathbf L\mathbf C-\mathbf B\mathbf K)\,]^{-1}\mathbf L.
  \label{eq:syn-comp-tf}
\end{equation}
$\mathbf C_K(s)$ 是一个\emph{真（proper）}的 $n$ 阶有理矩阵，即一个 $n$ 阶动态补偿器（\cref{def:syn-dyn-comp}）\cite{chen1999linear,astrom2008feedback}。\textbf{它自身的极点是 $\operatorname{eig}(\mathbf A-\mathbf L\mathbf C-\mathbf B\mathbf K)$}——注意这\emph{既不是}闭环极点、\emph{也不是}观测器极点。

\begin{pitfall}[把“补偿器自身极点”误当成“闭环极点”]
动态补偿器 $\mathbf C_K(s)$ 的极点是 $\operatorname{eig}(\mathbf A-\mathbf L\mathbf C-\mathbf B\mathbf K)$，而带它构成的闭环系统的极点是 $\operatorname{eig}(\mathbf A-\mathbf B\mathbf K)\cup\operatorname{eig}(\mathbf A-\mathbf L\mathbf C)$（分离原理，\cref{thm:syn-separation}）。两者\emph{一般不相等}！补偿器（孤立地看）甚至可以是\emph{不稳定}的（$\mathbf A-\mathbf L\mathbf C-\mathbf B\mathbf K$ 含右半平面特征值），而闭环照样稳定——这正是“用不稳定补偿器镇定稳定对象”在现代控制里的合法情形。设计时\emph{配的是闭环极点}（经 $\mathbf K,\mathbf L$），而非补偿器自身极点；后者是 $\mathbf K,\mathbf L$ 选定后的\emph{派生}产物。
\end{pitfall}

\begin{insight}
\textbf{观测器 $=$ 把静态升级为动态所需的那段“记忆”。}\;静态输出反馈 $\mathbf u=\mathbf v-\mathbf H\mathbf y$ 没有状态，故只能在低维的可达极点曲面上挪动；观测器恰好注入了 $n$ 维内部动态 $\hat{\mathbf x}$，把控制器变成 \cref{def:syn-dyn-comp} 意义下的动态补偿器，从而\emph{恢复}了全状态反馈的配极点能力。这就解释了 \cref{sec:syn-output-pole} 的悖论——“输出反馈配不动极点”说的是\emph{静态}输出反馈；\emph{动态}输出反馈（观测器）是另一个更大的类，能覆盖全空间。
\end{insight}

\paragraph{传递函数视角：拆成“串联补偿器 $+$ 反馈补偿器”（刘豹 §5.6.3）。}
\cref{eq:syn-comp-tf} 把观测器反馈打包成\emph{单个} $\mathbf C_K(s)$。刘豹则给出一个等价但更贴近经典框图的拆法：把它分解为一个置于\emph{前向通道}的串联补偿器与一个置于\emph{反馈通道}的反馈补偿器\cite{liubao2006modern}。把“带反馈阵 $\mathbf K$ 的观测器”看作子系统 $\Sigma_c$，其输入为 $(\mathbf u,\mathbf y)$。由 \cref{eq:syn-luenberger} 取拉氏变换（零初值）$\hat{\mathbf X}(s)=[\,s\mathbf I-(\mathbf A-\mathbf L\mathbf C)\,]^{-1}[\mathbf L\mathbf Y(s)+\mathbf B\mathbf U(s)]$，代入 $\mathbf U=\mathbf V-\mathbf K\hat{\mathbf X}$：
\begin{equation}
  \mathbf U(s)=\mathbf V(s)-\underbrace{\mathbf K[\,s\mathbf I-(\mathbf A-\mathbf L\mathbf C)\,]^{-1}\mathbf B}_{\textstyle \mathbf W_{\mathrm s}^{*}(s)}\,\mathbf U(s)-\underbrace{\mathbf K[\,s\mathbf I-(\mathbf A-\mathbf L\mathbf C)\,]^{-1}\mathbf L}_{\textstyle \mathbf W_{\mathrm f}(s)}\,\mathbf Y(s).
  \label{eq:syn-comp-collect}
\end{equation}
其中 $\mathbf W_{\mathrm f}(s)$ 是以 $\mathbf y$ 为输入的\textbf{反馈补偿器}。把含 $\mathbf U$ 的项归并：$[\mathbf I+\mathbf W_{\mathrm s}^{*}(s)]\mathbf U=\mathbf V-\mathbf W_{\mathrm f}(s)\mathbf Y$，于是
\begin{equation}
  \mathbf U(s)=\mathbf W_{\mathrm s}(s)\big[\mathbf V(s)-\mathbf W_{\mathrm f}(s)\mathbf Y(s)\big],\qquad \mathbf W_{\mathrm s}(s):=[\mathbf I+\mathbf W_{\mathrm s}^{*}(s)]^{-1},
  \label{eq:syn-comp-Ws-def}
\end{equation}
$\mathbf W_{\mathrm s}(s)$ 即置于前向通道的\textbf{串联补偿器}。\emph{这恰是一个带串联补偿器 $\mathbf W_{\mathrm s}$ 与反馈补偿器 $\mathbf W_{\mathrm f}$ 的输出反馈系统}（刘豹 图 5.24 / 5.25 结构）。

\paragraph{$\mathbf W_{\mathrm s}$ 物理可实现。}
$\mathbf W_{\mathrm s}=[\mathbf I+\mathbf W_{\mathrm s}^{*}]^{-1}$ 是否\emph{真有理}（可由有限维系统实现）？用矩阵求逆（push-through）恒等式可得其闭式：
\begin{equation}
  \boxed{\;\mathbf W_{\mathrm s}(s)=\mathbf I-\mathbf K\,[\,s\mathbf I-(\mathbf A-\mathbf L\mathbf C-\mathbf B\mathbf K)\,]^{-1}\mathbf B\;}
  \label{eq:syn-comp-Ws}
\end{equation}
它是\emph{常阵 $\mathbf I$ $+$ 严格真项}，故真有理、物理可实现，极点为 $\operatorname{eig}(\mathbf A-\mathbf L\mathbf C-\mathbf B\mathbf K)$（与单包补偿器 $\mathbf C_K$ 同极点）。\emph{验证}（记 $\mathbf N=s\mathbf I-(\mathbf A-\mathbf L\mathbf C)$，$\mathbf N'=s\mathbf I-(\mathbf A-\mathbf L\mathbf C-\mathbf B\mathbf K)=\mathbf N+\mathbf B\mathbf K$）：由 $\mathbf N^{-1}\mathbf B\mathbf K=\mathbf N^{-1}\mathbf N'-\mathbf I$ 右乘 ${\mathbf N'}^{-1}\mathbf B$ 得关键恒等式 $\mathbf K\mathbf N^{-1}\mathbf B\,\mathbf K{\mathbf N'}^{-1}\mathbf B=\mathbf K\mathbf N^{-1}\mathbf B-\mathbf K{\mathbf N'}^{-1}\mathbf B$，于是
\begin{align}
  [\mathbf I+\mathbf W_{\mathrm s}^{*}]\,\mathbf W_{\mathrm s}
  &=(\mathbf I+\mathbf K\mathbf N^{-1}\mathbf B)(\mathbf I-\mathbf K{\mathbf N'}^{-1}\mathbf B)\notag\\
  &=\mathbf I-\mathbf K{\mathbf N'}^{-1}\mathbf B+\mathbf K\mathbf N^{-1}\mathbf B-\mathbf K\mathbf N^{-1}\mathbf B\,\mathbf K{\mathbf N'}^{-1}\mathbf B=\mathbf I.
  \label{eq:syn-comp-verify}
\end{align}
即 \cref{eq:syn-comp-Ws} 确为 $[\mathbf I+\mathbf W_{\mathrm s}^{*}]^{-1}$（此即刘豹式 5.131--5.136 的可实现性证明）。两补偿器的状态空间实现（由 $\mathbf W_{\mathrm f},\mathbf W_{\mathrm s}$ 读出）为
\begin{equation}
  \underbrace{\dot{\mathbf z}_2=(\mathbf A-\mathbf L\mathbf C)\mathbf z_2+\mathbf L\,\mathbf y,\ \ \mathbf w_2=\mathbf K\mathbf z_2}_{\text{反馈补偿器 }\mathbf W_{\mathrm f}};\qquad
  \underbrace{\dot{\mathbf z}_1=(\mathbf A-\mathbf L\mathbf C-\mathbf B\mathbf K)\mathbf z_1+\mathbf B\,\boldsymbol\zeta,\ \ \mathbf w_1=-\mathbf K\mathbf z_1+\boldsymbol\zeta}_{\text{串联补偿器 }\mathbf W_{\mathrm s}},
  \label{eq:syn-comp-realize}
\end{equation}
其中串联补偿器输入 $\boldsymbol\zeta=\mathbf v-\mathbf w_2$、输出 $\mathbf w_1=\mathbf u$ 驱动对象。

\begin{theorem}[带观测器状态反馈 $\equiv$ 带补偿器输出反馈]
\label{thm:syn-compensator-equiv}
在传递（输入输出）特性意义下，带全维观测器的状态反馈系统\textbf{完全等效于}一个带串联补偿器 $\mathbf W_{\mathrm s}(s)$（\cref{eq:syn-comp-Ws}）与反馈补偿器 $\mathbf W_{\mathrm f}(s)$（\cref{eq:syn-comp-collect}）的输出反馈系统；反之，用这两个经典补偿器可\emph{构造}出与观测器反馈传递等效的系统\cite{liubao2006modern}。
\end{theorem}

\paragraph{状态空间级验证：传递函数仍是 $\mathbf W_K(s)$。}
把对象 $\Sigma_0$ 与两补偿器 \cref{eq:syn-comp-realize} 互联（$\boldsymbol\zeta=\mathbf v-\mathbf w_2$、$\mathbf u=\mathbf w_1$、反馈补偿器输入取 $\mathbf y$），以 $(\mathbf z_2,\mathbf x,\mathbf z_1)$ 为状态得 $3n$ 维闭环（记其系统阵为 $\mathbf A_{\mathrm{cl}}$）：
\begin{equation}
  \begin{bmatrix}\dot{\mathbf z}_2\\\dot{\mathbf x}\\\dot{\mathbf z}_1\end{bmatrix}
  =\underbrace{\begin{bmatrix}\mathbf A-\mathbf L\mathbf C & \mathbf L\mathbf C & \mathbf 0\\ -\mathbf B\mathbf K & \mathbf A & -\mathbf B\mathbf K\\ -\mathbf B\mathbf K & \mathbf 0 & \mathbf A-\mathbf L\mathbf C-\mathbf B\mathbf K\end{bmatrix}}_{\mathbf A_{\mathrm{cl}}}
  \begin{bmatrix}\mathbf z_2\\\mathbf x\\\mathbf z_1\end{bmatrix}+\begin{bmatrix}\mathbf 0\\\mathbf B\\\mathbf B\end{bmatrix}\mathbf v,\qquad \mathbf y=[\mathbf 0\ \ \mathbf C\ \ \mathbf 0]\begin{bmatrix}\mathbf z_2\\\mathbf x\\\mathbf z_1\end{bmatrix}.
  \label{eq:syn-comp-aug}
\end{equation}
取等效变换 $\mathbf T=\big[\begin{smallmatrix}\mathbf I&\mathbf 0&\mathbf 0\\\mathbf 0&\mathbf I&\mathbf 0\\-\mathbf I&\mathbf I&-\mathbf I\end{smallmatrix}\big]$（$\mathbf T^{-1}=\mathbf T$），对 $\mathbf A_{\mathrm{cl}}$ 作 $\bar{\mathbf A}_{\mathrm{cl}}=\mathbf T^{-1}\mathbf A_{\mathrm{cl}}\mathbf T$ 化为\emph{块上三角}：
\begin{equation}
  \bar{\mathbf A}_{\mathrm{cl}}=\begin{bmatrix}\mathbf A-\mathbf L\mathbf C & \mathbf L\mathbf C & \mathbf 0\\ \mathbf 0 & \mathbf A-\mathbf B\mathbf K & \mathbf B\mathbf K\\ \mathbf 0 & \mathbf 0 & \mathbf A-\mathbf L\mathbf C\end{bmatrix},\quad
  \bar{\mathbf B}=\mathbf T^{-1}\begin{bmatrix}\mathbf 0\\\mathbf B\\\mathbf B\end{bmatrix}=\begin{bmatrix}\mathbf 0\\\mathbf B\\\mathbf 0\end{bmatrix},\quad
  \bar{\mathbf C}=[\mathbf 0\ \mathbf C\ \mathbf 0]\,\mathbf T=[\mathbf 0\ \mathbf C\ \mathbf 0].
  \label{eq:syn-comp-tri}
\end{equation}
输入只进入中间块、输出只读取中间块，而中间对角块恰是 $\mathbf A-\mathbf B\mathbf K$，故
\begin{equation}
  \mathbf W_z(s)=\mathbf C[\,s\mathbf I-(\mathbf A-\mathbf B\mathbf K)\,]^{-1}\mathbf B=\mathbf W_K(s),
  \label{eq:syn-comp-WK}
\end{equation}
与直接状态反馈系统（\cref{eq:syn-sep-augmented} 的 $\mathbf v\!\to\!\mathbf y$ 通道）\emph{完全相同}。闭环极点 $=\operatorname{eig}(\mathbf A-\mathbf B\mathbf K)\cup\operatorname{eig}(\mathbf A-\mathbf L\mathbf C)\cup\operatorname{eig}(\mathbf A-\mathbf L\mathbf C)$：控制器极点 $+$ \emph{两份}观测器极点。这两份观测器极点对应不可控 / 不可观的隐藏模态，在 $\mathbf W_z$ 中被零点对消、不出现——与 \cref{thm:syn-separation} 处的“传递函数不变性”一脉相承。

\begin{insight}
\textbf{三种实现、三种维数、同一个 $\mathbf W_K(s)$。}\;(i) \emph{直接状态反馈}：$n$ 维，极点 $\operatorname{eig}(\mathbf A-\mathbf B\mathbf K)$，理想但需全状态；(ii) \emph{观测器反馈}：$2n$ 维（$\mathbf x$ 与 $\hat{\mathbf x}$），极点 $\operatorname{eig}(\mathbf A-\mathbf B\mathbf K)\cup\operatorname{eig}(\mathbf A-\mathbf L\mathbf C)$，控制器部分是 $n$ 阶补偿器 $\mathbf C_K$；(iii) \emph{串 $+$ 反补偿器}：$3n$ 维，极点再多一份 $\operatorname{eig}(\mathbf A-\mathbf L\mathbf C)$，是\emph{非最小}的经典实现。三者\emph{共享同一传递函数} $\mathbf W_K(s)$，却维数、可控可观性、最小性各异。\textbf{教训：传递函数等价 $\neq$ 状态空间等价}——I/O 行为相同的系统，内部结构（乃至鲁棒性、噪声响应）可以天差地别。这与 LQG 鲁棒裕度不可分离（\cref{sec:syn-robust} 的 Doyle 反例）是同一枚硬币：$\mathbf W_K(s)$ 看不到的隐藏模态，正是鲁棒性藏身之处。
\end{insight}

\begin{insight}
\textbf{与分离原理、与经典控制、与机器人级联架构的接口。}\;分离原理（\cref{thm:syn-separation}）讲的是\emph{闭环}极点可分离；本节讲的是\emph{控制器本身}是一个有自身极点 $\operatorname{eig}(\mathbf A-\mathbf L\mathbf C-\mathbf B\mathbf K)$ 的动态补偿器。向经典控制看：观测器型控制器是一类\emph{结构化的串 $+$ 反校正网络}（超前 / 滞后的多变量推广）；频域 $H_\infty$ / 回路成形（\cref{ch:robust-hinf}）设计一般的 $\mathbf C(s)$，观测器型只是其中一种参数化——而\emph{所有}镇定控制器都能写成“观测器 $+$ 状态反馈 $+$ Youla 参数 $\mathbf Q(s)$”的形式\cite{chen1999linear}。向机器人看：一切“状态估计器（观测器 / 卡尔曼，\cref{ch:eskf}）$+$ 反馈控制器”的\emph{级联}架构——如 SLAM 定位喂给轨迹跟踪——本质上\emph{就是}一个动态输出反馈补偿器，其 $n$ 维“记忆”即估计器状态。分离原理给了这种级联以合法性；但当鲁棒性吃紧时，必须把“估计 $+$ 控制”当作\emph{一个}动态系统统一设计（LQG / $H_\infty$，\cref{ch:lqr-lqg-riccati,ch:robust-hinf}）。
\end{insight}

\begin{example}[观测器型控制器的补偿器阶数]
\label{ex:syn-compensator}
仍取 \cref{ex:syn-full} 的对象 $W_0(s)=\dfrac1{s(s+6)}$（$n=2$，单输入单输出 $m=1$）。其\emph{全维}观测器型控制器是一个 $n=2$ 阶动态补偿器 $\mathbf C_K(s)=\mathbf K[s\mathbf I-(\mathbf A-\mathbf L\mathbf c-\mathbf b\mathbf K)]^{-1}\mathbf L$（\cref{eq:syn-comp-tf}）。而 \cref{thm:syn-dyn-comp-pole} 指出：\emph{任意}配置 $n=2$ 个极点其实只需 $n-1=1$ 阶动态补偿器。差额从何弥合？正是\emph{降维观测器}：\cref{ex:syn-full} 的降维观测器维数 $n-m=2-1=1$，对应的观测器型控制器恰是 $1$ 阶补偿器，\textbf{正好达到 \cref{thm:syn-dyn-comp-pole} 的最小阶 $n-1$}。可见全维观测器（$n$ 阶补偿器）在补偿器阶数上是\emph{冗余}的，降维观测器（$n-m$ 阶）才是配极点的\emph{最小阶}动态补偿器——这从“补偿器阶数”角度，给了 \cref{sec:syn-reduced} 降维观测器一个全新而更深的动机。
\end{example}

% ════════════════════════════════════════════════════════════════
\section{鲁棒入门：当模型本身不确定}
\label{sec:syn-robust}

\paragraph{动机：精确极点配置依赖精确模型——可模型从不精确。}
前面所有设计（极点配置、观测器、分离原理）都\emph{默认模型 $(\mathbf A,\mathbf B)$ 精确已知}。但实际建模总要简化、忽略高阶动态、且系统运行中受各种扰动——真实对象其实是一个\emph{不确定系统}\cite{zhang2017modern}：
\begin{equation}
  \dot{\mathbf x}=(\mathbf A+\Delta\mathbf A)\mathbf x+(\mathbf B+\Delta\mathbf B)\mathbf u,
  \label{eq:syn-uncertain}
\end{equation}
$\Delta\mathbf A,\Delta\mathbf B$ 是\emph{未知}的不确定项（有界，常用范数描述 $\|\Delta\mathbf A\|\le$ 某界）。此时无法对一个“不知道的 $\mathbf A$”精确配极点。\textbf{鲁棒镇定（robust stabilization）}问的是：能否设计一个固定的状态反馈，使闭环对\emph{整族}不确定性都渐近稳定？“鲁棒”（robust）即抗扰能力——鲁棒性越强，容忍的模型误差越大。本节是 P4 鲁棒主题的\emph{人审入口}，给出 Lyapunov / LMI / $H_\infty$ 三件工具的概念与机器人实例；频域 $H_\infty/\mu$ 综合的深推交 \cref{ch:robust-hinf}。

\paragraph{工具一：Schur 补——把非线性矩阵不等式线性化。}
鲁棒设计反复用到一个把“二次型矩阵不等式”化为“线性矩阵不等式”的杠杆。

\begin{lemma}[Schur 补]
\label{lem:syn-schur}
对称分块阵 $\mathbf M=\big[\begin{smallmatrix}\mathbf A_{11}&\mathbf A_{12}\\\mathbf A_{12}^\top&\mathbf A_{22}\end{smallmatrix}\big]\succ\mathbf 0$ 的充要条件为\cite{zhang2017modern,boyd1994linear}
\begin{equation}
  \mathbf A_{11}\succ\mathbf 0\ \text{且}\ \mathbf A_{22}-\mathbf A_{12}^\top\mathbf A_{11}^{-1}\mathbf A_{12}\succ\mathbf 0,\quad\text{或}\quad \mathbf A_{22}\succ\mathbf 0\ \text{且}\ \mathbf A_{11}-\mathbf A_{12}\mathbf A_{22}^{-1}\mathbf A_{12}^\top\succ\mathbf 0.
  \label{eq:syn-schur}
\end{equation}
其中 $\mathbf A_{22}-\mathbf A_{12}^\top\mathbf A_{11}^{-1}\mathbf A_{12}$ 称 $\mathbf A_{11}$ 在 $\mathbf M$ 中的\emph{Schur 补}。
\end{lemma}

\begin{insight}
\textbf{Schur 补为何是 LMI 的“瑞士军刀”？}\;一个含 $\mathbf P^{-1}$ 或 $\mathbf P\mathbf B\mathbf B^\top\mathbf P$ 这类\emph{对变量非线性}的不等式（如 Riccati 不等式），通过 Schur 补可\emph{等价}改写成一个对变量 $\mathbf P$（或 $\mathbf Q=\mathbf P^{-1}$）\emph{线性}的分块矩阵不等式——即 LMI。LMI 是\emph{凸}的，可用内点法（\cref{ch:nlopt}）高效全局求解（MATLAB LMI 工具箱、CVX）。这就是为什么现代鲁棒控制几乎全用 LMI 语言：把设计问题“凸化”，求解器一键搞定。
\end{insight}

\paragraph{工具二：Lyapunov + LMI 鲁棒镇定。}
对满足\emph{匹配条件}（$\Delta\mathbf A,\Delta\mathbf B$ 落在 $\mathbf B$ 的列空间内）的不确定系统，可用 Lyapunov 函数 $V(\mathbf x)=\mathbf x^\top\mathbf P\mathbf x$ 设计鲁棒反馈\cite{zhang2017modern}。核心步骤：先解一个名义 Riccati 方程
\begin{equation}
  \mathbf P\mathbf A+\mathbf A^\top\mathbf P-\mathbf P\mathbf B\mathbf B^\top\mathbf P+\mathbf Q_0=\mathbf 0
  \label{eq:syn-robust-riccati}
\end{equation}
\footnote{源（现控02 式 6.6.6）OCR 把 $\mathbf B\mathbf B^\top$ 误印为 $\mathbf B^\top\mathbf B$（M2-c L3506）：$\mathbf B$ 为 $n\times p$ 时 $\mathbf B^\top\mathbf B$ 维度不符，且下文增益 $\mathbf K=-\mathbf B^\top\mathbf P$ 与证明中出现 $\mathbf P\mathbf B\mathbf B^\top\mathbf P$ 均佐证应为 $\mathbf B\mathbf B^\top$，本书据勘误改正。}
得正定解 $\mathbf P\succ\mathbf 0$，取 $\mathbf K=-\mathbf B^\top\mathbf P$\footnote{源此处 OCR 把增益误写成 $\mathbf K=-\mathbf B^\top\mathbf P\mathbf x$（混入状态 $\mathbf x$，M2-c L3509）；增益是\emph{矩阵}、不含 $\mathbf x$，本书改正。}（使名义闭环 $\mathbf A_1=\mathbf A+\mathbf B\mathbf K$ 稳定）。考察 $V(\mathbf x)=\mathbf x^\top\mathbf P\mathbf x$ 沿不确定系统 \cref{eq:syn-uncertain} 的导数，把交叉项用引理
\begin{equation}
  \mathbf X^\top\mathbf Y+\mathbf Y^\top\mathbf X\preceq\tfrac1r\mathbf X^\top\mathbf X+r\mathbf Y^\top\mathbf Y\quad(\forall r>0)
  \label{eq:syn-young}
\end{equation}
逐一“配方上界”，可证存在足够大的 $\gamma$ 使
\begin{equation}
  \dot V\le\mathbf x^\top\big[\mathbf P\mathbf A_1+\mathbf A_1^\top\mathbf P+\mathbf P\mathbf B\mathbf R^{-1}\mathbf B^\top\mathbf P+\mathbf Q\big]\mathbf x<0
  \label{eq:syn-robust-vdot}
\end{equation}
对一切容许的 $\Delta\mathbf A,\Delta\mathbf B$ 成立（只要选 $\mathbf Q$ 盖住不确定界、$\gamma\ge(2-\xi)/\delta$、$\mathbf R^{-1}=\mathbf I/\xi$ 且第二个 Riccati 方程 $\mathbf A_1^\top\mathbf P+\mathbf P\mathbf A_1+\mathbf P\mathbf B\mathbf R^{-1}\mathbf B^\top\mathbf P+\mathbf Q=\mathbf 0$ 有正定解\cite{zhang2017modern}）。于是得鲁棒反馈 $\mathbf u=(\mathbf K-\gamma\mathbf B^\top\mathbf P)\mathbf x$，闭环在原点\emph{大范围}渐近稳定。对\emph{不满足}匹配条件的情形，可把 $\Delta\mathbf A=\mathbf E_1\boldsymbol\Sigma_1\mathbf F_1,\ \Delta\mathbf B=\mathbf E_2\boldsymbol\Sigma_2\mathbf F_2$（$\boldsymbol\Sigma_i^\top\boldsymbol\Sigma_i\preceq\mathbf I$）代入，化为另一个含 $\gamma_i$ 的 Riccati / LMI 可行性问题（反馈律 $\mathbf u=-\mathbf B^\top\mathbf P\mathbf x$）；至于含\emph{状态时滞}的鲁棒镇定，亦循同一“李氏 $\to$ Schur 补 $\to$ LMI”套路，单列于 \cref{sec:syn-delay}。\emph{无论何种描述，用 Lyapunov 设计的思路一致}：造 $V$、令 $\dot V<0$ 鲁棒成立、化 LMI。这条“Lyapunov $\to$ Riccati 不等式 $\to$ Schur 补 $\to$ LMI”的流水线，是把镇定从“需精确模型”升级为“容忍不确定”的核心套路（接 \cref{ch:lyapunov-basics} 的李氏方法、\cref{ch:nlopt} 的凸优化）。

\subsection{时滞系统的鲁棒镇定（LMI）}
\label{sec:syn-delay}

\paragraph{动机：状态里带“过去”。}
许多实际系统的演化不仅取决于\emph{此刻}状态，还取决于\emph{一段时间之前}的状态——传输延迟、测量与执行滞后、网络通信时延都会引入\textbf{状态时滞}（如蒸汽锅炉温控、网络化机器人）。带不确定性的线性时滞系统写作\cite{zhang2017modern}
\begin{equation}
  \dot{\mathbf x}(t)=(\mathbf A+\Delta\mathbf A)\mathbf x(t)+\mathbf A_h\,\mathbf x(t-h)+\mathbf B\mathbf u,
  \label{eq:syn-delay-sys}
\end{equation}
其中 $h>0$ 为时滞常数，$\mathbf A_h$ 为时滞耦合阵（满足范数界 $\mathbf A_h^\top\mathbf A_h\preceq\mathbf I$），不确定项仍取范数有界结构 $\Delta\mathbf A=\mathbf E_1\boldsymbol\Sigma_1\mathbf F_1$（$\boldsymbol\Sigma_1^\top\boldsymbol\Sigma_1\preceq\mathbf I$，与上文一致）。时滞项 $\mathbf x(t-h)$ 使系统成为\emph{无穷维}的泛函微分方程，单纯的 $V(\mathbf x)=\mathbf x^\top\mathbf P\mathbf x$ 不够用——需要一个能“记住过去一段轨迹”的李氏泛函。

\paragraph{工具：Lyapunov--Krasovskii 泛函。}
取\textbf{Lyapunov--Krasovskii 泛函}（在二次型上叠加一段历史积分）
\begin{equation}
  V(\mathbf x_t)=\mathbf x^\top(t)\mathbf P\mathbf x(t)+\int_{t-h}^{t}\mathbf x^\top(\theta)\mathbf x(\theta)\,\mathrm d\theta,\qquad \mathbf P\succ\mathbf 0,
  \label{eq:syn-delay-lk}
\end{equation}
其积分项正是用来\emph{补偿}时滞 $\mathbf x(t-h)$ 的：沿 \cref{eq:syn-delay-sys}（先令 $\mathbf u=\mathbf 0$）求导，积分项贡献 $\mathbf x^\top(t)\mathbf x(t)-\mathbf x^\top(t-h)\mathbf x(t-h)$，后一块恰好“吃掉”交叉项里 $\mathbf x(t-h)$ 的二次式。

\begin{theorem}[时滞系统鲁棒稳定的 LMI 判据（张嗣瀛\cite{zhang2017modern} §6.6.3）]
\label{thm:syn-delay-stab}
对 $\mathbf u=\mathbf 0$ 的时滞系统 \cref{eq:syn-delay-sys}，若存在 $\mathbf P\succ\mathbf 0$ 与正数 $\varepsilon_1,\varepsilon_2$ 使
\begin{equation}
  \begin{bmatrix}\mathbf A^\top\mathbf P+\mathbf P\mathbf A+\mathbf I+\varepsilon_2\mathbf F_1^\top\mathbf F_1 & \mathbf P & \mathbf P\mathbf E_1\\ \mathbf P & -\varepsilon_1\mathbf I & \mathbf 0\\ \mathbf E_1^\top\mathbf P & \mathbf 0 & -\varepsilon_2\mathbf I\end{bmatrix}\prec\mathbf 0,
  \label{eq:syn-delay-lmi-analysis}
\end{equation}
则系统在原点（对一切容许的 $\Delta\mathbf A,\mathbf A_h$）渐近稳定。
\end{theorem}
\begin{proof}[要点（LK $+$ 配方 $+$ Schur 补）]
沿 \cref{eq:syn-delay-lk} 求导并代入 \cref{eq:syn-delay-sys}（$\mathbf u=\mathbf 0$）。用 Young 不等式 \cref{eq:syn-young} 把含 $\mathbf x(t-h)$ 的交叉项上界为 $\varepsilon_1^{-1}\mathbf x^\top\mathbf P^2\mathbf x+\varepsilon_1\mathbf x^\top(t-h)\mathbf A_h^\top\mathbf A_h\mathbf x(t-h)$，再用 $\mathbf A_h^\top\mathbf A_h\preceq\mathbf I$ 与积分项的 $-\mathbf x^\top(t-h)\mathbf x(t-h)$ 相消（取 $\varepsilon_1\le1$ 即令 $(\varepsilon_1-1)\le0$ 项可弃），并把不确定项 $\Delta\mathbf A=\mathbf E_1\boldsymbol\Sigma_1\mathbf F_1$ 也配方上界，得
\[
  \dot V\le\mathbf x^\top(t)\big(\mathbf A^\top\mathbf P+\mathbf P\mathbf A+\mathbf I+\varepsilon_1^{-1}\mathbf P^2+\varepsilon_2\mathbf F_1^\top\mathbf F_1+\varepsilon_2^{-1}\mathbf P\mathbf E_1\mathbf E_1^\top\mathbf P\big)\mathbf x(t).
\]
括号 $\prec\mathbf 0$ 即 $\dot V\prec0$。对其中 $\varepsilon_1^{-1}\mathbf P^2$ 与 $\varepsilon_2^{-1}\mathbf P\mathbf E_1\mathbf E_1^\top\mathbf P$ 两个二次项各做一次 Schur 补（\cref{lem:syn-schur}），即等价升维为 \cref{eq:syn-delay-lmi-analysis} 的 $3\times3$ 块 LMI。\hfill$\square$
\end{proof}

\begin{theorem}[时滞系统状态反馈镇定的 LMI（张嗣瀛\cite{zhang2017modern} §6.6.3）]
\label{thm:syn-delay-synth}
对时滞系统 \cref{eq:syn-delay-sys} 施状态反馈 $\mathbf u=\mathbf K\mathbf x$，闭环 $\dot{\mathbf x}=(\mathbf A+\Delta\mathbf A+\mathbf B\mathbf K)\mathbf x+\mathbf A_h\mathbf x(t-h)$。对给定 $\varepsilon_1,\varepsilon_2>0$，若存在 $\mathbf Q\succ\mathbf 0$ 与 $\mathbf Y$ 使
\begin{equation}
  \begin{bmatrix}\mathbf Q\mathbf A^\top+\mathbf Y^\top\mathbf B^\top+\mathbf A\mathbf Q+\mathbf B\mathbf Y & \mathbf I & \mathbf E_1 & \mathbf Q\mathbf F_1^\top & \mathbf Q\\ \mathbf I & -\varepsilon_1\mathbf I & \mathbf 0 & \mathbf 0 & \mathbf 0\\ \mathbf E_1^\top & \mathbf 0 & -\varepsilon_2\mathbf I & \mathbf 0 & \mathbf 0\\ \mathbf F_1\mathbf Q & \mathbf 0 & \mathbf 0 & -\varepsilon_2^{-1}\mathbf I & \mathbf 0\\ \mathbf Q & \mathbf 0 & \mathbf 0 & \mathbf 0 & -\mathbf I\end{bmatrix}\prec\mathbf 0,
  \label{eq:syn-delay-lmi-synth}
\end{equation}
则系统可镇定，反馈增益 $\mathbf K=\mathbf Y\mathbf Q^{-1}$。
\end{theorem}
\begin{proof}[要点（合同换元线性化双线性项）]
把 \cref{thm:syn-delay-stab} 中的 $\mathbf A$ 换成闭环 $\mathbf A_1=\mathbf A+\mathbf B\mathbf K$，得到一个对 $(\mathbf P,\mathbf K)$ \emph{双线性}（含 $\mathbf P\mathbf B\mathbf K$）的不等式——非凸、不能直接 LMI 求解。\textbf{关键一步}：对该不等式左右合同变换（左右乘 $\operatorname{diag}(\mathbf P^{-1},\mathbf I,\dots)$），令 $\mathbf Q=\mathbf P^{-1}$、$\mathbf Y=\mathbf K\mathbf Q$，则 $\mathbf P^{-1}\mathbf A_1^\top+\mathbf A_1\mathbf P^{-1}=\mathbf Q\mathbf A^\top+\mathbf Y^\top\mathbf B^\top+\mathbf A\mathbf Q+\mathbf B\mathbf Y$——双线性项被\emph{线性化}为对 $(\mathbf Q,\mathbf Y)$ 的线性式，配合对 $\mathbf P^2$ 等项的 Schur 补展开即 \cref{eq:syn-delay-lmi-synth}。MATLAB LMI 工具箱解出 $\mathbf Q,\mathbf Y$ 后 $\mathbf K=\mathbf Y\mathbf Q^{-1}$。\hfill$\square$
\end{proof}

\begin{insight}
\textbf{“变量代换 $\mathbf Q=\mathbf P^{-1},\,\mathbf Y=\mathbf K\mathbf Q$”是 LMI 综合的通用解锁手法.}\;镇定综合天然出现 $\mathbf P\mathbf B\mathbf K$ 这类\emph{决策变量相乘}的项（非凸，求解器无能为力）。换元到 $\mathbf Q=\mathbf P^{-1}$、把 $\mathbf K$ 吸收进 $\mathbf Y=\mathbf K\mathbf Q$，可把一大类“Lyapunov 镇定 / 鲁棒镇定 / 时滞镇定 / $H_\infty$ 综合”统一化为标准 LMI 可行性问题——这与 \cref{lem:syn-schur} 的 Schur 补一起，正是现代鲁棒控制能“一键”交给凸求解器的根本原因。时滞只是把判据从代数 LMI 多加了一段 LK 历史积分，套路不变。
\end{insight}

\paragraph{工具三：$H_\infty$ 状态反馈。}
当不确定性 / 干扰用\emph{频域}范数刻画时，进入 $H_\infty$ 控制：设广义对象有干扰输入 $\mathbf w$、评价输出 $\mathbf z$、传递阵 $\mathbf T_{zw}(s)$，目标是设计反馈使闭环稳定\emph{且} $\|\mathbf T_{zw}\|_\infty<1$（最坏频率下的最大增益受界，即抗扰能力的最优保证）\cite{zhang2017modern}。在适当假设下，状态反馈 $H_\infty$ 解由一个\emph{代数 Riccati 方程}给出：

\begin{theorem}[$H_\infty$ 状态反馈（标准型）]
\label{thm:syn-hinf}
设广义对象 $(\mathbf A,\mathbf B_1,\mathbf B_2,\mathbf C_1)$ 满足标准假设（$(\mathbf A,\mathbf B_2)$ 可镇定等），则存在状态反馈 $\mathbf u=-\mathbf K\mathbf x$ 使闭环稳定且 $\|\mathbf T_{zw}\|_\infty<1$ 的充要条件，是代数 Riccati 方程
\begin{equation}
  \mathbf A^\top\mathbf P+\mathbf P\mathbf A+\mathbf P(\mathbf B_1\mathbf B_1^\top-\mathbf B_2\mathbf B_2^\top)\mathbf P+\mathbf C_1^\top\mathbf C_1=\mathbf 0
  \label{eq:syn-hinf-are}
\end{equation}
\footnote{源（现控02 式 6.6.27 邻近）OCR 漏掉闭括号后的 $\mathbf P$（M2-c L4029）；标准 $H_\infty$ 状态反馈 ARE 第三项须为 $\mathbf P(\cdots)\mathbf P$，且下文出现 $\mathbf A+(\mathbf B_1\mathbf B_1^\top-\mathbf B_2\mathbf B_2^\top)\mathbf P$ 佐证，本书据勘误补全。}
有使 $\mathbf A+(\mathbf B_1\mathbf B_1^\top-\mathbf B_2\mathbf B_2^\top)\mathbf P$ 稳定的半正定解 $\mathbf P\succeq\mathbf 0$，此时控制器 $\mathbf u=-\mathbf B_2^\top\mathbf P\mathbf x$\cite{zhang2017modern}。
\end{theorem}

\begin{note}
\textbf{$H_\infty$ Riccati 与 LQR Riccati 的关系.}\;\cref{eq:syn-hinf-are} 与 LQR 的 CARE（导论 \cref{eq:ctrl-care}）形似，但多了一项 $+\mathbf P\mathbf B_1\mathbf B_1^\top\mathbf P$（来自干扰通道 $\mathbf B_1$）。这个“$+$”号正是\emph{最坏干扰}的痕迹：$H_\infty$ 控制可视为控制器与“对抗性干扰”之间的\emph{零和微分博弈}，干扰想最大化 $\|\mathbf z\|$、控制想最小化它，鞍点解给出博弈型 Riccati。LQR 是无干扰（$\mathbf B_1=\mathbf 0$）的退化情形。这一“博弈 Riccati”视角的严格展开见 \cref{ch:lqr-lqg-riccati,ch:robust-hinf}。
\end{note}

\begin{example}[回旋倒立摆 / 两足机器人的 $H_\infty$ 状态反馈（张嗣瀛\cite{zhang2017modern} §6.6.4 完整算例）]
\label{ex:syn-pendulum}
倒立摆是\textbf{两足行走机器人、火箭垂直姿态控制}等的最简模型\cite{zhang2017modern}。回旋型倒立摆是单输入双输出系统：输入是驱动直臂的电机力矩 $\tau$，输出是直臂转角 $\theta_0$ 与摆角 $\theta_1$。其动力学正是标准的\emph{机械臂形式}（与 \cref{ch:rigid_body} 一致）：
\begin{equation}
  \mathbf M(\boldsymbol\theta)\ddot{\boldsymbol\theta}+\mathbf C(\boldsymbol\theta,\dot{\boldsymbol\theta})\dot{\boldsymbol\theta}+\mathbf D\dot{\boldsymbol\theta}+\mathbf g(\boldsymbol\theta)=\boldsymbol\tau,
  \label{eq:syn-pendulum-dyn}
\end{equation}
$\boldsymbol\theta=[\theta_0\ \theta_1]^\top$，$\mathbf M(\boldsymbol\theta)=\big[\begin{smallmatrix}j_1+j_2\sin^2\theta_1&j_3\cos\theta_1\\j_3\cos\theta_1&j_2+j_4\end{smallmatrix}\big]$（惯量阵），$\mathbf D=\operatorname{diag}(d_0,d_1)$ 为阻尼，$\mathbf g(\boldsymbol\theta)=[0\ -j_5\sin\theta_1]^\top$ 为重力项，参数 $j_1=J_0+m_0l_0^2+m_1L_0^2$、$j_2=m_1l_1^2$、$j_3=m_1L_0l_1$、$j_4=J_1$、$j_5=m_1gl_1$。控制目标：摆保持倒立、直臂转角 $\theta_0=0$，即令平衡点 $\boldsymbol\theta=\dot{\boldsymbol\theta}=\mathbf 0$ 渐近稳定。

\emph{第一步（线性化 $\to$ 状态方程）.}\;在倒立平衡点处取 $\sin\theta_1\approx\theta_1,\cos\theta_1\approx1$ 线性化，状态 $\mathbf x=[\theta_0,\theta_1,\dot\theta_0,\dot\theta_1]^\top$、把作用于摆上的扰动力矩 $\omega$（模型近似误差的等价扰动）记为干扰输入，得
\begin{equation}
  \dot{\mathbf x}=\mathbf A\mathbf x+\mathbf B_1\omega+\mathbf B_2 u,\quad
  \mathbf A=\frac1a\!\begin{bmatrix}0&0&a&0\\0&0&0&a\\0&-j_3j_5&-(j_2{+}j_4)d_0&j_3d_1\\0&j_1j_5&j_3d_0&-j_1d_1\end{bmatrix}\!,\ \
  \mathbf B_1=\!\begin{bmatrix}0\\0\\-j_3\\j_1\end{bmatrix}\!,\ \
  \mathbf B_2=\frac1a\!\begin{bmatrix}0\\0\\j_2{+}j_4\\-j_3\end{bmatrix}\!,
  \label{eq:syn-pendulum-ss}
\end{equation}
其中 $a=j_1(j_2+j_4)-j_3^2$，$u=\tau$。倒立平衡点处 $\mathbf A$ 含\emph{正}特征值（开环不稳，摆会倒）。

\emph{第二步（设计指标 $\to$ 评价信号）.}\;要求 (i) $\mathbf x=\mathbf 0$ 局部渐近稳定；(ii) 对任意 $\omega\in L_2[0,\infty)$ 有扰动抑制
\[
  \int_0^\infty\!\big(q_1\theta_0^2+q_2\theta_1^2+q_3\dot\theta_0^2+q_4\dot\theta_1^2+\rho u^2\big)\,\mathrm dt<\int_0^\infty\!\omega^2\,\mathrm dt\quad(q_i\ge0,\ \rho>0\ \text{为加权}).
\]
定义评价信号 $\mathbf z=\mathbf C_1\mathbf x+\mathbf D_{12}u$，取
\[
  \mathbf C_1=\begin{bmatrix}\sqrt{q_1}&0&0&0\\0&\sqrt{q_2}&0&0\\0&0&\sqrt{q_3}&0\\0&0&0&\sqrt{q_4}\\0&0&0&0\end{bmatrix},\qquad
  \mathbf D_{12}=[\,0,\,0,\,0,\,0,\,\sqrt\rho\,]^\top,
\]
则 $\mathbf C_1^\top\mathbf C_1=\operatorname{diag}(q_1,q_2,q_3,q_4)$、$\mathbf D_{12}^\top\mathbf D_{12}=\rho$、$\mathbf D_{12}^\top\mathbf C_1=\mathbf 0$（恰满足 \cref{thm:syn-hinf} 的标准假设 $\mathbf D_{12}^\top[\mathbf C_1\ \mathbf D_{12}]=[\mathbf 0\ \mathbf I]$），指标 (ii) 化为 $\|\mathbf z\|_2<\|\omega\|_2$，即 $\|\mathbf T_{zw}\|_\infty<1$。于是求满足 (i)(ii) 的 $\mathbf K$ 等价于求使闭环稳定且 $\|\mathbf T_{zw}\|_\infty<1$ 的 $H_\infty$ 状态反馈。

\emph{第三步（解 ARE $\to$ 数值增益）.}\;按系统辨识取 $j_1=0.1400,\ j_2=0.0061,\ j_3=0.0067,\ j_4=0.0001,\ j_5=0.2292,\ d_0=0.0464,\ d_1=0.0013$，加权 $q_1=2.25,\ q_2=4900,\ q_3=1,\ q_4=14400,\ \rho=1.44$。代入 \cref{eq:syn-pendulum-ss} 得数值 $\mathbf A,\mathbf B_1,\mathbf B_2$，数值求解 \cref{thm:syn-hinf} 的 ARE \cref{eq:syn-hinf-are} 的镇定半正定解 $\mathbf P\succeq\mathbf 0$，得反馈增益
\begin{equation}
  \mathbf K=\mathbf B_2^\top\mathbf P=[\,6.817\quad 665.0\quad 19.20\quad 162.6\,],\qquad u=-\mathbf K\mathbf x,
  \label{eq:syn-pendulum-K}
\end{equation}
即把开环不稳的倒立摆\emph{鲁棒镇定}于倒立位、且对惯量参数不准 / 忽略摩擦造成的等价扰动 $\omega$ 有 $H_\infty$ 抑制保证。$\mathbf K$ 中 $\theta_1$ 一列增益（$665$）最大，正反映“摆角偏离最危险、须最重地纠正”这一物理直觉。这正是“现代控制理论基础”一路通到\emph{平衡 / 足式机器人控制}的桥梁：同一套状态空间 $+$ Riccati 工具，在 \cref{eq:syn-pendulum-dyn} 这类机械臂动力学上实例化（频域 DGKF 双 Riccati 解与 $\mu$ 综合的深推交 \cref{ch:robust-hinf}）。
\end{example}

% ════════════════════════════════════════════════════════════════
\section{通向高级：本章在综合版图中的位置}
\label{sec:syn-outlook}

本章把线性系统综合“讲透”：从全状态反馈的极点配置出发，经镇定、解耦，到观测器（动机链：状态不可测$\to$微分器放噪$\to$开环漂移$\to$龙伯格纠错$\to$降维），再用分离原理把控制器与观测器焊成可实现的输出反馈，最后以鲁棒入门处理模型不确定。这层地基之上，A--Q 高级章把每件工具严格化或推广到前沿：

\begin{itemize}
  \item \textbf{分离原理 / LQG} $\to$ \cref{ch:lqr-lqg-riccati}：本章 \cref{thm:syn-separation} 的“极点可分离”在那里被严格化为\emph{确定性等价}（certainty equivalence）的完整证明——把最优状态反馈（LQR）与最优观测器（卡尔曼 / LQE）拼成 LQG，并证其最优性；同时给出 CARE 存在唯一性（可镇定 + 可检测 $\Rightarrow$ 唯一正定镇定解，本章 \cref{thm:syn-stabilize} 的“可镇定”是其前提）。
  \item \textbf{$H_\infty$ / 鲁棒} $\to$ \cref{ch:robust-hinf}：本章 \cref{sec:syn-robust} 的 LMI / $H_\infty$ 状态反馈（\cref{eq:syn-hinf-are}）在那里展开为\emph{频域} $H_\infty$ 综合、$\mu$ 分析与 DGKF 双 Riccati 解；本章一句带过的“LQG 鲁棒裕度不可分离”（Doyle 反例）在那里获得完整剖析。
  \item \textbf{观测器} $\to$ \cref{ch:eskf}：本章的龙伯格观测器与卡尔曼滤波\emph{结构同形}（预测 + 残差纠正），差别仅在增益来路——龙伯格由极点配置\emph{手选} $\mathbf L$，卡尔曼由噪声协方差经 Riccati \emph{最优}给出 $\mathbf L$。换言之：\textbf{卡尔曼滤波 = 最优的龙伯格观测器}，正如 LQR = 最优的极点配置。这是“估计--控制对偶”在结构层的最直接体现。
  \item \textbf{安全对偶} $\to$ \cref{ch:clf-cbf}：本章用 Lyapunov 函数证镇定 / 鲁棒（\cref{eq:syn-robust-riccati} 邻近）的思路，在那里对偶到\emph{控制屏障函数}（CBF）——Lyapunov 管“稳定”（吸引到平衡），CBF 管“安全”（不出安全集），二者是同一“在反馈下令某标量函数单调”范式的两面。
\end{itemize}

\noindent 在机器人学语境，本章工具是后续应用章的理论上游：极点配置 / 观测器是足式 WBC 与平衡控制的经典基线（与 \cref{ch:control} 足式部分的 LQR/QP 互补），$H_\infty$ / 鲁棒镇定是面对参数不确定（负载变化、地形未知）的保底手段，分离原理则是一切“状态估计器 + 反馈控制器”级联架构（SLAM 定位 + 轨迹跟踪）的合法性依据。读完本章，你应能在“可控 $\to$ 配极点、可观 $\to$ 造观测器、分离原理 $\to$ 拼装、不确定 $\to$ 鲁棒化”这条链上自由游走，并知道每个结论该到哪一章求严格证明或前沿推广。

% ════════════════════════════════════════════════════════════════
\section{练习}
\label{sec:syn-exercises}

\begin{practice}
\begin{enumerate}
  \item （极点配置）对双积分器 $\mathbf A=\big[\begin{smallmatrix}0&1\\0&0\end{smallmatrix}\big],\ \mathbf b=\big[\begin{smallmatrix}0\\1\end{smallmatrix}\big]$（机器人“位置 + 速度”最简模型），用状态反馈把闭环极点配到 $-2\pm j2$。分别用可控标准型系数比对法（\cref{eq:syn-kbar}）与 Ackermann 公式（\cref{eq:syn-ackermann}）求 $\mathbf K$，验证两法一致。
  \item （镇定 vs 配置）给一个不可控但可镇定的系统（如 $\mathbf A=\operatorname{diag}(1,-1),\ \mathbf b=[1\ 0]^\top$），说明为何不能任意配极点、却仍可镇定；指出哪个特征值被“锁死”（\cref{thm:syn-stabilize}）。
  \item （对偶）对第 1 题取 $\mathbf c=[1\ 0]$（只测位置），设计全维观测器使误差极点为 $-10,-10$。先写对偶系统 $(\mathbf A^\top,\mathbf c^\top)$，对它配极点求 $\mathbf L^\top$，转置得 $\mathbf L$（\cref{thm:syn-obs-design}）。
  \item （降维）承上，$\operatorname{rank}\mathbf c=1$，故 $\bar{\mathbf x}_2$ 可由 $y$ 直接得、只需 $1$ 维降维观测器。按 \cref{sec:syn-reduced} 两步法设计，并说明它比全维观测器省在哪、对噪声的代价在哪。
  \item （分离原理）把第 1、3 题的 $\mathbf K,\mathbf L$ 拼成带观测器反馈，写出增广闭环阵 \cref{eq:syn-sep-augmented}，验证其特征值确为 $\{-2\pm j2\}\cup\{-10,-10\}$。再解释“观测器要比控制器快”这条整定法则的依据。
  \item （解耦）给一个 $\mathbf E$ 非奇异的双输入双输出系统，按 \cref{thm:syn-decouple} 求 $(\mathbf K,\mathbf F)$ 化为积分型，再对各子系统追加极点配置；另给一个 $\mathbf E$ 奇异的系统，说明它为何不能状态反馈解耦、该改用什么办法。
  \item （鲁棒，思考）对倒立摆线性化模型（\cref{ex:syn-pendulum}），把惯量参数 $j_i$ 的 $\pm10\%$ 误差写成 $\Delta\mathbf A$，讨论精确极点配置与 $H_\infty$ 状态反馈（\cref{thm:syn-hinf}）在该不确定下的表现差异，并指出 Schur 补在把镇定条件化为 LMI 时所起的作用。
\end{enumerate}
\end{practice}

% ════════════════════════════════════════════════════════════════
\section*{本章小结}

\begin{table}[H]\centering\small
\caption{符号表}
\begin{tabular}{ll}
\toprule
符号 & 含义 \\
\midrule
$\Sigma_0=(\mathbf A,\mathbf B,\mathbf C)$ & 受控对象（单输入单输出记 $(\mathbf A,\mathbf b,\mathbf c)$） \\
$\mathbf K$ & 状态反馈增益（闭环 $\mathbf A-\mathbf B\mathbf K$，律 $\mathbf u=\mathbf v-\mathbf K\mathbf x$） \\
$\mathbf L$ & 观测器增益（误差阵 $\mathbf A-\mathbf L\mathbf C$） \\
$\mathbf v$ & 参考输入 \\
$\mathbf H$ & 输出反馈增益（$\mathbf u=\mathbf v-\mathbf H\mathbf y$） \\
$\hat{\mathbf x},\ \tilde{\mathbf x}=\mathbf x-\hat{\mathbf x}$ & 状态估计 / 估计误差 \\
$\lambda_i^*,\ f^*(\lambda)$ & 期望闭环极点 / 期望特征多项式 \\
$\mathbf T_c,\ \mathcal C,\ \mathcal O$ & 可控标准型变换阵 / 可控性矩阵 / 可观性矩阵 \\
$d_i,\ \mathbf E,\ \mathbf F$ & 解耦阶数 / 解耦矩阵 / 输入变换阵 \\
$\mathbf P\succ\mathbf 0,\ \mathbf Q$ & 李氏 / Riccati 矩阵；李氏方程 $\mathbf A^\top\mathbf P+\mathbf P\mathbf A=-\mathbf Q$ \\
$\Delta\mathbf A,\Delta\mathbf B$ & 模型不确定项 \\
$\mathbf T_{zw}(s),\ \|\cdot\|_\infty$ & 干扰到评价输出的闭环传递阵 / $H_\infty$ 范数 \\
\bottomrule
\end{tabular}
\end{table}

\begin{table}[H]\centering\small
\caption{定理与判据速查表}
\begin{tabular}{lll}
\toprule
名称 & 标签 & 一句话 \\
\midrule
极点配置定理 & \cref{thm:syn-pole-placement} & 可任意配极点 $\iff$ 完全可控 \\
Ackermann 公式 & \cref{eq:syn-ackermann} & $\mathbf K=[0\cdots0\ 1]\mathcal C^{-1}f^*(\mathbf A)$ \\
状态反馈可镇定 & \cref{thm:syn-stabilize} & $\iff$ 不可控子系统渐近稳定 \\
可解耦判据 & \cref{thm:syn-decouple} & $\iff$ 解耦矩阵 $\mathbf E$ 非奇异 \\
观测器误差动态 & \cref{eq:syn-obs-error} & $\dot{\tilde{\mathbf x}}=(\mathbf A-\mathbf L\mathbf C)\tilde{\mathbf x}$，不依赖输入 \\
观测器增益设计 & \cref{thm:syn-obs-design} & 对偶极点配置（$(\mathbf A^\top,\mathbf C^\top)$ 可控） \\
观测器存在条件 & \cref{thm:syn-obs-exist} & $\iff$ 不可观子系统渐近稳定 \\
分离原理 & \cref{thm:syn-separation} & 闭环极点 = 控制器极点 $\cup$ 观测器极点 \\
动态补偿器配极点 & \cref{thm:syn-dyn-comp-pole} & 输出反馈经 $n-1$ 阶动态补偿器可任意配极点 \\
观测器即动态补偿器 & \cref{thm:syn-compensator-equiv} & 带观测器状态反馈 $\equiv$ 带串 / 反补偿器输出反馈，同 $\mathbf W_K(s)$ \\
Schur 补 & \cref{lem:syn-schur} & 把二次型 MI 化为 LMI \\
时滞镇定 LMI & \cref{thm:syn-delay-synth} & LK 泛函 $+$ 换元 $\mathbf Q{=}\mathbf P^{-1},\mathbf Y{=}\mathbf K\mathbf Q$，$\mathbf K{=}\mathbf Y\mathbf Q^{-1}$ \\
$H_\infty$ 状态反馈 & \cref{thm:syn-hinf} & 博弈型 ARE \cref{eq:syn-hinf-are}，$\mathbf u=-\mathbf B_2^\top\mathbf P\mathbf x$ \\
\bottomrule
\end{tabular}
\end{table}

\begin{table}[H]\centering\small
\caption{设计工具适用边界（知识点总表）}
\begin{tabular}{llll}
\toprule
工具 & 前提 & 给出 & 局限 \\
\midrule
极点配置 & 完全可控 & 任意指定闭环极点 & 需全状态、需精确模型 \\
镇定 & 不可控子系统稳 & 仅渐近稳定 & 不精确控制极点位置 \\
状态反馈解耦 & $\mathbf E$ 非奇异 & 独立单回路 & 条件苛刻、解耦后须再配极点 \\
全维观测器 & $(\mathbf A,\mathbf C)$ 可观 & 重构全状态 & $n$ 维、对噪声经模型平滑 \\
降维观测器 & 同上、$\operatorname{rank}\mathbf C=m$ & 重构 $n-m$ 维 & 低维省算，但 $\mathbf y$ 噪声直入 \\
分离原理 & 可控且可观（LTI） & $\mathbf K,\mathbf L$ 独立设计 & 鲁棒裕度不可分离、仅 LTI 严格 \\
$H_\infty$/LMI 鲁棒 & 可镇定 + 不确定有界 & 对整族不确定稳定 & 保守、需 $\Delta$ 的范数界 \\
\bottomrule
\end{tabular}
\end{table}

\section*{常见误解汇总}
\begin{enumerate}
  \item \textbf{“可控就能任意快地把状态搬到任意点。”}\; 错。可控只保证\emph{有限时间能到}，配得越快控制量越大、越敏感（\cref{sec:syn-ackermann} 陷阱）。
  \item \textbf{“测不到的状态可以微分输出得到。”}\; 错。理论可行但微分放大噪声、毫无工程价值，这正是观测器存在的理由（\cref{sec:syn-observer} 反面一）。
  \item \textbf{“观测器增益 $\mathbf L$ 要和控制器增益 $\mathbf K$ 一起算。”}\; 错。分离原理保证二者\emph{独立}设计（\cref{thm:syn-separation}）。
  \item \textbf{“分离原理成立 $\Rightarrow$ LQG 和 LQR 一样鲁棒。”}\; 错。极点可分离，但\emph{鲁棒裕度不可分离}（Doyle 反例，\cref{sec:syn-robust} 与 \cref{ch:robust-hinf}）。
  \item \textbf{“解耦做完就行了。”}\; 错。积分型解耦系统极点全在原点、性能差，须对各子系统追加极点配置（\cref{thm:syn-decouple} 后注）。
  \item \textbf{“静态输出反馈也能镇定任意可控可观系统。”}\; 错。静态输出反馈镇定能力弱，存在可控可观却无法静态输出反馈镇定的系统（\cref{sec:syn-stabilize}）。但\emph{动态}输出反馈（即观测器型动态补偿器）可以——它靠注入内部状态突破静态反馈的极点配置极限（\cref{thm:syn-dyn-comp-pole,sec:syn-compensator}）。
  \item \textbf{“观测器型控制器的极点就是闭环极点。”}\; 错。补偿器自身极点是 $\operatorname{eig}(\mathbf A-\mathbf L\mathbf C-\mathbf B\mathbf K)$，闭环极点才是控制器极点 $\cup$ 观测器极点；补偿器单看甚至可不稳定（\cref{sec:syn-compensator} 陷阱）。
\end{enumerate}

\section*{延伸阅读}
\begin{itemize}
  \item \textbf{主源（物理直觉 + 动机链 + 例题密集）}：刘豹《现代控制理论》第3版\cite{liubao2006modern} 第5章——状态反馈基本结构、极点配置定理与证明、镇定、解耦、观测器动机链（微分器放噪$\to$龙伯格）、全 / 降维观测器、分离原理的完整中文叙事，本章正文主线的出处。
  \item \textbf{第二信源（互校 + Ackermann + 鲁棒入门）}：张嗣瀛 / 高立群《现代控制理论》第2版\cite{zhang2017modern} 第6章——Ackermann / Bass--Gura 算法、反馈解耦 $\mathbf E$ 阵非奇异充要条件、降维观测器，以及 6.6 节 Schur 补 / LMI / $H_\infty$ 状态反馈 / 回旋倒立摆与两足机器人实例（本章鲁棒一节的人审 ground-truth）。
  \item \textbf{经典溯源}：Luenberger 观测器原典\cite{luenberger1971introduction}；Wonham 极点配置定理\cite{wonham1967pole}；Ackermann 公式\cite{ackermann1972der}；Boyd 等《Linear Matrix Inequalities in System and Control Theory》\cite{boyd1994linear}（LMI 方法奠基）。
  \item \textbf{现代教材}：Åström \& Murray《Feedback Systems》\cite{astrom2008feedback}（状态反馈 / 观测器 / 分离原理的现代讲法）；Anderson \& Moore《Optimal Control》\cite{anderson2007optimal}（LQG 与分离原理）；Khalil《Nonlinear Systems》\cite{khalil2002nonlinear}（Lyapunov 鲁棒设计）。
\end{itemize}

\section*{本章与后续章节的关系}
\begin{table}[H]\centering\small
\begin{tabular}{lll}
\toprule
后续章节 & 关系 & 本章铺垫 \\
\midrule
\cref{ch:lqr-lqg-riccati}（LQR/LQG/Riccati）& 严格化 & 分离原理 \cref{thm:syn-separation}、可镇定 \cref{thm:syn-stabilize} \\
\cref{ch:robust-hinf}（鲁棒 / $H_\infty$）& 严格化 & LMI / $H_\infty$ 状态反馈 \cref{eq:syn-hinf-are} \\
\cref{ch:clf-cbf}（CLF/CBF）& 对偶 & Lyapunov 镇定 \cref{eq:syn-robust-riccati} 邻近 \\
\cref{ch:eskf}（状态估计）& 横向对偶 & 观测器 \cref{eq:syn-luenberger}（卡尔曼 = 最优龙伯格） \\
\cref{ch:optimal-control-basics}（最优控制）& 互补 & 极点配置（手选极点）$\to$ LQR（最优选极点） \\
\cref{ch:control}（控制导论）& 展开 & 极点配置 / 观测器 / 分离原理（\cref{sec:ctrl-pole-placement}） \\
\bottomrule
\end{tabular}
\end{table}

\section*{故障排查手册}
\begin{table}[H]\centering\small
\begin{tabular}{p{0.36\linewidth}p{0.58\linewidth}}
\toprule
症状 & 可能原因与对策 \\
\midrule
极点配不到指定位置 & 系统不完全可控：不可控子系统极点锁死（\cref{thm:syn-pole-placement} 必要性 insight）。先做结构分解看 $\tilde{\mathbf A}_{22}$。 \\
闭环可控但变得不可观 & 状态反馈改极点不改零点，制造了零极点对消（\cref{sec:syn-feedback} 末）。检查闭环传递函数是否约分。 \\
观测器估计被噪声淹没 & 误差极点配得太快、$\mathbf L$ 过大把测量噪声灌进 $\hat{\mathbf x}$（\cref{sec:syn-observer} 陷阱）。放慢误差极点，或改卡尔曼滤波按噪声协方差最优折中（\cref{ch:eskf}）。 \\
带观测器闭环不稳 & 检查是否 $(\mathbf A,\mathbf C)$ 不可观（观测器不存在，\cref{thm:syn-obs-exist}）或 $\mathbf L$ 未使 $\mathbf A-\mathbf L\mathbf C$ 稳定；分离原理要求可控\emph{且}可观。 \\
解耦后系统性能极差 / 临界不稳 & 积分型解耦极点全在原点，须追加极点配置（\cref{thm:syn-decouple} 后注）；且子系统须可控。 \\
解耦设计失败 & 解耦矩阵 $\mathbf E$ 奇异 $\Rightarrow$ 不能状态反馈解耦（\cref{thm:syn-decouple}）。改用前馈补偿器，或前馈 + 状态反馈混合。 \\
模型误差导致闭环不稳 & 精确极点配置不鲁棒。改鲁棒镇定：解 Riccati \cref{eq:syn-robust-riccati} / $H_\infty$ ARE \cref{eq:syn-hinf-are}，或写成 LMI 求解（\cref{sec:syn-robust}）。 \\
LMI 求解器报不可行 & 不确定界过大 / 不满足匹配条件。收紧 $\Delta$ 的范数界，或换不满足匹配条件的设计路线（\cref{sec:syn-robust}）。 \\
\bottomrule
\end{tabular}
\end{table}
